
% Default to the notebook output style

    


% Inherit from the specified cell style.




\documentclass[11pt,letter,landscape]{article}    
%\documentclass[11pt]{article}

    
    
    \usepackage[T1]{fontenc}
    % Nicer default font (+ math font) than Computer Modern for most use cases
    \usepackage{mathpazo}

    % Basic figure setup, for now with no caption control since it's done
    % automatically by Pandoc (which extracts ![](path) syntax from Markdown).
    \usepackage{graphicx}
    % We will generate all images so they have a width \maxwidth. This means
    % that they will get their normal width if they fit onto the page, but
    % are scaled down if they would overflow the margins.
    \makeatletter
    \def\maxwidth{\ifdim\Gin@nat@width>\linewidth\linewidth
    \else\Gin@nat@width\fi}
    \makeatother
    \let\Oldincludegraphics\includegraphics
    % Set max figure width to be 80% of text width, for now hardcoded.
    \renewcommand{\includegraphics}[1]{\Oldincludegraphics[width=.8\maxwidth]{#1}}
    % Ensure that by default, figures have no caption (until we provide a
    % proper Figure object with a Caption API and a way to capture that
    % in the conversion process - todo).
    \usepackage{caption}
    \DeclareCaptionLabelFormat{nolabel}{}
    \captionsetup{labelformat=nolabel}

    \usepackage{adjustbox} % Used to constrain images to a maximum size 
    \usepackage{xcolor} % Allow colors to be defined
    \usepackage{enumerate} % Needed for markdown enumerations to work
    \usepackage{geometry} % Used to adjust the document margins
    \usepackage{amsmath} % Equations
    \usepackage{amssymb} % Equations
    \usepackage{textcomp} % defines textquotesingle
    % Hack from http://tex.stackexchange.com/a/47451/13684:
    \AtBeginDocument{%
        \def\PYZsq{\textquotesingle}% Upright quotes in Pygmentized code
    }
    \usepackage{upquote} % Upright quotes for verbatim code
    \usepackage{eurosym} % defines \euro
    \usepackage[mathletters]{ucs} % Extended unicode (utf-8) support
    \usepackage[utf8x]{inputenc} % Allow utf-8 characters in the tex document
    \usepackage{fancyvrb} % verbatim replacement that allows latex
    \usepackage{grffile} % extends the file name processing of package graphics 
                         % to support a larger range 
    % The hyperref package gives us a pdf with properly built
    % internal navigation ('pdf bookmarks' for the table of contents,
    % internal cross-reference links, web links for URLs, etc.)
    \usepackage{hyperref}
    \usepackage{longtable} % longtable support required by pandoc >1.10
    \usepackage{booktabs}  % table support for pandoc > 1.12.2
    \usepackage[inline]{enumitem} % IRkernel/repr support (it uses the enumerate* environment)
    \usepackage[normalem]{ulem} % ulem is needed to support strikethroughs (\sout)
                                % normalem makes italics be italics, not underlines
    

    
    
    % Colors for the hyperref package
    \definecolor{urlcolor}{rgb}{0,.145,.698}
    \definecolor{linkcolor}{rgb}{.71,0.21,0.01}
    \definecolor{citecolor}{rgb}{.12,.54,.11}

    % ANSI colors
    \definecolor{ansi-black}{HTML}{3E424D}
    \definecolor{ansi-black-intense}{HTML}{282C36}
    \definecolor{ansi-red}{HTML}{E75C58}
    \definecolor{ansi-red-intense}{HTML}{B22B31}
    \definecolor{ansi-green}{HTML}{00A250}
    \definecolor{ansi-green-intense}{HTML}{007427}
    \definecolor{ansi-yellow}{HTML}{DDB62B}
    \definecolor{ansi-yellow-intense}{HTML}{B27D12}
    \definecolor{ansi-blue}{HTML}{208FFB}
    \definecolor{ansi-blue-intense}{HTML}{0065CA}
    \definecolor{ansi-magenta}{HTML}{D160C4}
    \definecolor{ansi-magenta-intense}{HTML}{A03196}
    \definecolor{ansi-cyan}{HTML}{60C6C8}
    \definecolor{ansi-cyan-intense}{HTML}{258F8F}
    \definecolor{ansi-white}{HTML}{C5C1B4}
    \definecolor{ansi-white-intense}{HTML}{A1A6B2}

    % commands and environments needed by pandoc snippets
    % extracted from the output of `pandoc -s`
    \providecommand{\tightlist}{%
      \setlength{\itemsep}{0pt}\setlength{\parskip}{0pt}}
    \DefineVerbatimEnvironment{Highlighting}{Verbatim}{commandchars=\\\{\}}
    % Add ',fontsize=\small' for more characters per line
    \newenvironment{Shaded}{}{}
    \newcommand{\KeywordTok}[1]{\textcolor[rgb]{0.00,0.44,0.13}{\textbf{{#1}}}}
    \newcommand{\DataTypeTok}[1]{\textcolor[rgb]{0.56,0.13,0.00}{{#1}}}
    \newcommand{\DecValTok}[1]{\textcolor[rgb]{0.25,0.63,0.44}{{#1}}}
    \newcommand{\BaseNTok}[1]{\textcolor[rgb]{0.25,0.63,0.44}{{#1}}}
    \newcommand{\FloatTok}[1]{\textcolor[rgb]{0.25,0.63,0.44}{{#1}}}
    \newcommand{\CharTok}[1]{\textcolor[rgb]{0.25,0.44,0.63}{{#1}}}
    \newcommand{\StringTok}[1]{\textcolor[rgb]{0.25,0.44,0.63}{{#1}}}
    \newcommand{\CommentTok}[1]{\textcolor[rgb]{0.38,0.63,0.69}{\textit{{#1}}}}
    \newcommand{\OtherTok}[1]{\textcolor[rgb]{0.00,0.44,0.13}{{#1}}}
    \newcommand{\AlertTok}[1]{\textcolor[rgb]{1.00,0.00,0.00}{\textbf{{#1}}}}
    \newcommand{\FunctionTok}[1]{\textcolor[rgb]{0.02,0.16,0.49}{{#1}}}
    \newcommand{\RegionMarkerTok}[1]{{#1}}
    \newcommand{\ErrorTok}[1]{\textcolor[rgb]{1.00,0.00,0.00}{\textbf{{#1}}}}
    \newcommand{\NormalTok}[1]{{#1}}
    
    % Additional commands for more recent versions of Pandoc
    \newcommand{\ConstantTok}[1]{\textcolor[rgb]{0.53,0.00,0.00}{{#1}}}
    \newcommand{\SpecialCharTok}[1]{\textcolor[rgb]{0.25,0.44,0.63}{{#1}}}
    \newcommand{\VerbatimStringTok}[1]{\textcolor[rgb]{0.25,0.44,0.63}{{#1}}}
    \newcommand{\SpecialStringTok}[1]{\textcolor[rgb]{0.73,0.40,0.53}{{#1}}}
    \newcommand{\ImportTok}[1]{{#1}}
    \newcommand{\DocumentationTok}[1]{\textcolor[rgb]{0.73,0.13,0.13}{\textit{{#1}}}}
    \newcommand{\AnnotationTok}[1]{\textcolor[rgb]{0.38,0.63,0.69}{\textbf{\textit{{#1}}}}}
    \newcommand{\CommentVarTok}[1]{\textcolor[rgb]{0.38,0.63,0.69}{\textbf{\textit{{#1}}}}}
    \newcommand{\VariableTok}[1]{\textcolor[rgb]{0.10,0.09,0.49}{{#1}}}
    \newcommand{\ControlFlowTok}[1]{\textcolor[rgb]{0.00,0.44,0.13}{\textbf{{#1}}}}
    \newcommand{\OperatorTok}[1]{\textcolor[rgb]{0.40,0.40,0.40}{{#1}}}
    \newcommand{\BuiltInTok}[1]{{#1}}
    \newcommand{\ExtensionTok}[1]{{#1}}
    \newcommand{\PreprocessorTok}[1]{\textcolor[rgb]{0.74,0.48,0.00}{{#1}}}
    \newcommand{\AttributeTok}[1]{\textcolor[rgb]{0.49,0.56,0.16}{{#1}}}
    \newcommand{\InformationTok}[1]{\textcolor[rgb]{0.38,0.63,0.69}{\textbf{\textit{{#1}}}}}
    \newcommand{\WarningTok}[1]{\textcolor[rgb]{0.38,0.63,0.69}{\textbf{\textit{{#1}}}}}
    
    
    % Define a nice break command that doesn't care if a line doesn't already
    % exist.
    \def\br{\hspace*{\fill} \\* }
    % Math Jax compatability definitions
    \def\gt{>}
    \def\lt{<}
    % Document parameters
    \title{JOD Best Practices Lab}
    
    
    

    % Pygments definitions
    
\makeatletter
\def\PY@reset{\let\PY@it=\relax \let\PY@bf=\relax%
    \let\PY@ul=\relax \let\PY@tc=\relax%
    \let\PY@bc=\relax \let\PY@ff=\relax}
\def\PY@tok#1{\csname PY@tok@#1\endcsname}
\def\PY@toks#1+{\ifx\relax#1\empty\else%
    \PY@tok{#1}\expandafter\PY@toks\fi}
\def\PY@do#1{\PY@bc{\PY@tc{\PY@ul{%
    \PY@it{\PY@bf{\PY@ff{#1}}}}}}}
\def\PY#1#2{\PY@reset\PY@toks#1+\relax+\PY@do{#2}}

\expandafter\def\csname PY@tok@w\endcsname{\def\PY@tc##1{\textcolor[rgb]{0.73,0.73,0.73}{##1}}}
\expandafter\def\csname PY@tok@c\endcsname{\let\PY@it=\textit\def\PY@tc##1{\textcolor[rgb]{0.25,0.50,0.50}{##1}}}
\expandafter\def\csname PY@tok@cp\endcsname{\def\PY@tc##1{\textcolor[rgb]{0.74,0.48,0.00}{##1}}}
\expandafter\def\csname PY@tok@k\endcsname{\let\PY@bf=\textbf\def\PY@tc##1{\textcolor[rgb]{0.00,0.50,0.00}{##1}}}
\expandafter\def\csname PY@tok@kp\endcsname{\def\PY@tc##1{\textcolor[rgb]{0.00,0.50,0.00}{##1}}}
\expandafter\def\csname PY@tok@kt\endcsname{\def\PY@tc##1{\textcolor[rgb]{0.69,0.00,0.25}{##1}}}
\expandafter\def\csname PY@tok@o\endcsname{\def\PY@tc##1{\textcolor[rgb]{0.40,0.40,0.40}{##1}}}
\expandafter\def\csname PY@tok@ow\endcsname{\let\PY@bf=\textbf\def\PY@tc##1{\textcolor[rgb]{0.67,0.13,1.00}{##1}}}
\expandafter\def\csname PY@tok@nb\endcsname{\def\PY@tc##1{\textcolor[rgb]{0.00,0.50,0.00}{##1}}}
\expandafter\def\csname PY@tok@nf\endcsname{\def\PY@tc##1{\textcolor[rgb]{0.00,0.00,1.00}{##1}}}
\expandafter\def\csname PY@tok@nc\endcsname{\let\PY@bf=\textbf\def\PY@tc##1{\textcolor[rgb]{0.00,0.00,1.00}{##1}}}
\expandafter\def\csname PY@tok@nn\endcsname{\let\PY@bf=\textbf\def\PY@tc##1{\textcolor[rgb]{0.00,0.00,1.00}{##1}}}
\expandafter\def\csname PY@tok@ne\endcsname{\let\PY@bf=\textbf\def\PY@tc##1{\textcolor[rgb]{0.82,0.25,0.23}{##1}}}
\expandafter\def\csname PY@tok@nv\endcsname{\def\PY@tc##1{\textcolor[rgb]{0.10,0.09,0.49}{##1}}}
\expandafter\def\csname PY@tok@no\endcsname{\def\PY@tc##1{\textcolor[rgb]{0.53,0.00,0.00}{##1}}}
\expandafter\def\csname PY@tok@nl\endcsname{\def\PY@tc##1{\textcolor[rgb]{0.63,0.63,0.00}{##1}}}
\expandafter\def\csname PY@tok@ni\endcsname{\let\PY@bf=\textbf\def\PY@tc##1{\textcolor[rgb]{0.60,0.60,0.60}{##1}}}
\expandafter\def\csname PY@tok@na\endcsname{\def\PY@tc##1{\textcolor[rgb]{0.49,0.56,0.16}{##1}}}
\expandafter\def\csname PY@tok@nt\endcsname{\let\PY@bf=\textbf\def\PY@tc##1{\textcolor[rgb]{0.00,0.50,0.00}{##1}}}
\expandafter\def\csname PY@tok@nd\endcsname{\def\PY@tc##1{\textcolor[rgb]{0.67,0.13,1.00}{##1}}}
\expandafter\def\csname PY@tok@s\endcsname{\def\PY@tc##1{\textcolor[rgb]{0.73,0.13,0.13}{##1}}}
\expandafter\def\csname PY@tok@sd\endcsname{\let\PY@it=\textit\def\PY@tc##1{\textcolor[rgb]{0.73,0.13,0.13}{##1}}}
\expandafter\def\csname PY@tok@si\endcsname{\let\PY@bf=\textbf\def\PY@tc##1{\textcolor[rgb]{0.73,0.40,0.53}{##1}}}
\expandafter\def\csname PY@tok@se\endcsname{\let\PY@bf=\textbf\def\PY@tc##1{\textcolor[rgb]{0.73,0.40,0.13}{##1}}}
\expandafter\def\csname PY@tok@sr\endcsname{\def\PY@tc##1{\textcolor[rgb]{0.73,0.40,0.53}{##1}}}
\expandafter\def\csname PY@tok@ss\endcsname{\def\PY@tc##1{\textcolor[rgb]{0.10,0.09,0.49}{##1}}}
\expandafter\def\csname PY@tok@sx\endcsname{\def\PY@tc##1{\textcolor[rgb]{0.00,0.50,0.00}{##1}}}
\expandafter\def\csname PY@tok@m\endcsname{\def\PY@tc##1{\textcolor[rgb]{0.40,0.40,0.40}{##1}}}
\expandafter\def\csname PY@tok@gh\endcsname{\let\PY@bf=\textbf\def\PY@tc##1{\textcolor[rgb]{0.00,0.00,0.50}{##1}}}
\expandafter\def\csname PY@tok@gu\endcsname{\let\PY@bf=\textbf\def\PY@tc##1{\textcolor[rgb]{0.50,0.00,0.50}{##1}}}
\expandafter\def\csname PY@tok@gd\endcsname{\def\PY@tc##1{\textcolor[rgb]{0.63,0.00,0.00}{##1}}}
\expandafter\def\csname PY@tok@gi\endcsname{\def\PY@tc##1{\textcolor[rgb]{0.00,0.63,0.00}{##1}}}
\expandafter\def\csname PY@tok@gr\endcsname{\def\PY@tc##1{\textcolor[rgb]{1.00,0.00,0.00}{##1}}}
\expandafter\def\csname PY@tok@ge\endcsname{\let\PY@it=\textit}
\expandafter\def\csname PY@tok@gs\endcsname{\let\PY@bf=\textbf}
\expandafter\def\csname PY@tok@gp\endcsname{\let\PY@bf=\textbf\def\PY@tc##1{\textcolor[rgb]{0.00,0.00,0.50}{##1}}}
\expandafter\def\csname PY@tok@go\endcsname{\def\PY@tc##1{\textcolor[rgb]{0.53,0.53,0.53}{##1}}}
\expandafter\def\csname PY@tok@gt\endcsname{\def\PY@tc##1{\textcolor[rgb]{0.00,0.27,0.87}{##1}}}
\expandafter\def\csname PY@tok@err\endcsname{\def\PY@bc##1{\setlength{\fboxsep}{0pt}\fcolorbox[rgb]{1.00,0.00,0.00}{1,1,1}{\strut ##1}}}
\expandafter\def\csname PY@tok@kc\endcsname{\let\PY@bf=\textbf\def\PY@tc##1{\textcolor[rgb]{0.00,0.50,0.00}{##1}}}
\expandafter\def\csname PY@tok@kd\endcsname{\let\PY@bf=\textbf\def\PY@tc##1{\textcolor[rgb]{0.00,0.50,0.00}{##1}}}
\expandafter\def\csname PY@tok@kn\endcsname{\let\PY@bf=\textbf\def\PY@tc##1{\textcolor[rgb]{0.00,0.50,0.00}{##1}}}
\expandafter\def\csname PY@tok@kr\endcsname{\let\PY@bf=\textbf\def\PY@tc##1{\textcolor[rgb]{0.00,0.50,0.00}{##1}}}
\expandafter\def\csname PY@tok@bp\endcsname{\def\PY@tc##1{\textcolor[rgb]{0.00,0.50,0.00}{##1}}}
\expandafter\def\csname PY@tok@fm\endcsname{\def\PY@tc##1{\textcolor[rgb]{0.00,0.00,1.00}{##1}}}
\expandafter\def\csname PY@tok@vc\endcsname{\def\PY@tc##1{\textcolor[rgb]{0.10,0.09,0.49}{##1}}}
\expandafter\def\csname PY@tok@vg\endcsname{\def\PY@tc##1{\textcolor[rgb]{0.10,0.09,0.49}{##1}}}
\expandafter\def\csname PY@tok@vi\endcsname{\def\PY@tc##1{\textcolor[rgb]{0.10,0.09,0.49}{##1}}}
\expandafter\def\csname PY@tok@vm\endcsname{\def\PY@tc##1{\textcolor[rgb]{0.10,0.09,0.49}{##1}}}
\expandafter\def\csname PY@tok@sa\endcsname{\def\PY@tc##1{\textcolor[rgb]{0.73,0.13,0.13}{##1}}}
\expandafter\def\csname PY@tok@sb\endcsname{\def\PY@tc##1{\textcolor[rgb]{0.73,0.13,0.13}{##1}}}
\expandafter\def\csname PY@tok@sc\endcsname{\def\PY@tc##1{\textcolor[rgb]{0.73,0.13,0.13}{##1}}}
\expandafter\def\csname PY@tok@dl\endcsname{\def\PY@tc##1{\textcolor[rgb]{0.73,0.13,0.13}{##1}}}
\expandafter\def\csname PY@tok@s2\endcsname{\def\PY@tc##1{\textcolor[rgb]{0.73,0.13,0.13}{##1}}}
\expandafter\def\csname PY@tok@sh\endcsname{\def\PY@tc##1{\textcolor[rgb]{0.73,0.13,0.13}{##1}}}
\expandafter\def\csname PY@tok@s1\endcsname{\def\PY@tc##1{\textcolor[rgb]{0.73,0.13,0.13}{##1}}}
\expandafter\def\csname PY@tok@mb\endcsname{\def\PY@tc##1{\textcolor[rgb]{0.40,0.40,0.40}{##1}}}
\expandafter\def\csname PY@tok@mf\endcsname{\def\PY@tc##1{\textcolor[rgb]{0.40,0.40,0.40}{##1}}}
\expandafter\def\csname PY@tok@mh\endcsname{\def\PY@tc##1{\textcolor[rgb]{0.40,0.40,0.40}{##1}}}
\expandafter\def\csname PY@tok@mi\endcsname{\def\PY@tc##1{\textcolor[rgb]{0.40,0.40,0.40}{##1}}}
\expandafter\def\csname PY@tok@il\endcsname{\def\PY@tc##1{\textcolor[rgb]{0.40,0.40,0.40}{##1}}}
\expandafter\def\csname PY@tok@mo\endcsname{\def\PY@tc##1{\textcolor[rgb]{0.40,0.40,0.40}{##1}}}
\expandafter\def\csname PY@tok@ch\endcsname{\let\PY@it=\textit\def\PY@tc##1{\textcolor[rgb]{0.25,0.50,0.50}{##1}}}
\expandafter\def\csname PY@tok@cm\endcsname{\let\PY@it=\textit\def\PY@tc##1{\textcolor[rgb]{0.25,0.50,0.50}{##1}}}
\expandafter\def\csname PY@tok@cpf\endcsname{\let\PY@it=\textit\def\PY@tc##1{\textcolor[rgb]{0.25,0.50,0.50}{##1}}}
\expandafter\def\csname PY@tok@c1\endcsname{\let\PY@it=\textit\def\PY@tc##1{\textcolor[rgb]{0.25,0.50,0.50}{##1}}}
\expandafter\def\csname PY@tok@cs\endcsname{\let\PY@it=\textit\def\PY@tc##1{\textcolor[rgb]{0.25,0.50,0.50}{##1}}}

\def\PYZbs{\char`\\}
\def\PYZus{\char`\_}
\def\PYZob{\char`\{}
\def\PYZcb{\char`\}}
\def\PYZca{\char`\^}
\def\PYZam{\char`\&}
\def\PYZlt{\char`\<}
\def\PYZgt{\char`\>}
\def\PYZsh{\char`\#}
\def\PYZpc{\char`\%}
\def\PYZdl{\char`\$}
\def\PYZhy{\char`\-}
\def\PYZsq{\char`\'}
\def\PYZdq{\char`\"}
\def\PYZti{\char`\~}
% for compatibility with earlier versions
\def\PYZat{@}
\def\PYZlb{[}
\def\PYZrb{]}
\makeatother


    % Exact colors from NB
    \definecolor{incolor}{rgb}{0.0, 0.0, 0.5}
    \definecolor{outcolor}{rgb}{0.545, 0.0, 0.0}



    
    % Prevent overflowing lines due to hard-to-break entities
    \sloppy 
    % Setup hyperref package
    \hypersetup{
      breaklinks=true,  % so long urls are correctly broken across lines
      colorlinks=true,
      urlcolor=urlcolor,
      linkcolor=linkcolor,
      citecolor=citecolor,
      }
    % Slightly bigger margins than the latex defaults
    
    \geometry{verbose,tmargin=1in,bmargin=1in,lmargin=1in,rmargin=1in}
    
    

    \begin{document}
    
    
    \maketitle
    
    

    
    \section{JOD Best Practices Lab}\label{jod-best-practices-lab}

\begin{figure}
\centering
\includegraphics{inclusions/jodteenytinycube.png}
\caption{}
\end{figure}

    \subsubsection{Introduction}\label{introduction}

Software tools are like loaded guns: powerful weapons for slaying your
coding demons but also dangerous when used improperly. Have you ever
shot yourself in the foot? I know I'm missing a few toes and I bet you
are as well.

This lab outlines a number of "best practices" or guidelines for using
JOD. I've learned the hard way; if you take my advice to heart you
\emph{might} be spared!

    \begin{Verbatim}[commandchars=\\\{\}]
{\color{incolor}In [{\color{incolor}1}]:} \PY{c+c1}{NB. show J version}
        \PY{l+m+mi}{9}\PY{o}{!}\PY{o}{:}\PY{l+m+mi}{14} \PY{l+s}{\PYZsq{}}\PY{l+s}{\PYZsq{}}
\end{Verbatim}


    \begin{Verbatim}[commandchars=\\\{\}]
j901/j64avx/windows/beta-s/commercial/www.jsoftware.com/2019-12-02T12:51:33

    \end{Verbatim}

    \begin{Verbatim}[commandchars=\\\{\}]
{\color{incolor}In [{\color{incolor}2}]:} \PY{n+nv}{load} \PY{l+s}{\PYZsq{}}\PY{l+s}{g}\PY{l+s}{e}\PY{l+s}{n}\PY{l+s}{e}\PY{l+s}{r}\PY{l+s}{a}\PY{l+s}{l}\PY{l+s}{/}\PY{l+s}{j}\PY{l+s}{o}\PY{l+s}{d}\PY{l+s}{\PYZsq{}}
        
        \PY{c+c1}{NB. use portable box drawing characters}
        \PY{n+nv}{portchars} \PY{l+s}{\PYZsq{}}\PY{l+s}{\PYZsq{}}
\end{Verbatim}


    \begin{Verbatim}[commandchars=\\\{\}]


    \end{Verbatim}

    \subsubsection{Test for JOD user defined
folders}\label{test-for-jod-user-defined-folders}

\textbf{WARNING: THIS LAB REQUIRES A NUMBER OF USER DEFINED JOD
FOLDERS.}

This lab requires a number of JOD folders to run. The next step tests
for the existence of these folders.

If (\texttt{TestJODDirectories}) lists any undefined JOD directories
configure these directories before running this lab.
\href{https://github.com/jsoftware/general_joddocument/blob/master/pdfdoc/jod.pdf}{Instructions
on how to do this} can be found in (\texttt{jod.pdf}). Install the
\href{https://code.jsoftware.com/wiki/Addons/general/joddocument}{joddocument}
addon to get (\texttt{jod.pdf}).

    \begin{Verbatim}[commandchars=\\\{\}]
{\color{incolor}In [{\color{incolor}3}]:} \PY{n+nv}{TestJODDirectories\PYZus{}ijod\PYZus{}}\PY{o}{=:}\PY{n+nf}{3 : 0}
        
        \PY{c+c1}{NB.*TestJODDirectories v\PYZhy{}\PYZhy{} test user configured JOD directories.}
        \PY{c+c1}{NB.}
        \PY{c+c1}{NB. This  verb  checks  that  required  JOD  lab directories  are}
        \PY{c+c1}{NB. defined. \PYZdq{}Defined\PYZdq{} does not mean the  directories  exist only}
        \PY{c+c1}{NB. that (jpath) expands to something other than its default.}
        \PY{c+c1}{NB.}
        \PY{c+c1}{NB. monad:  clMsg =. TestJODDirectories uuIgnore}
        
        \PY{c+c1}{NB. when a relative directory does not exist (jpath) echoes its argument}
        \PY{n+nv}{jodudirs} \PY{o}{=.} \PY{l+s}{\PYZsq{}}\PY{l+s}{\PYZti{}}\PY{l+s}{\PYZsq{}}\PY{o}{\PYZam{}}\PY{o}{,}\PY{o}{\PYZam{}}\PY{o}{.}\PY{o}{\PYZgt{}} \PY{o}{;}\PY{o}{:}\PY{l+s}{\PYZsq{}}\PY{l+s}{J}\PY{l+s}{O}\PY{l+s}{D}\PY{l+s}{ }\PY{l+s}{J}\PY{l+s}{O}\PY{l+s}{D}\PY{l+s}{D}\PY{l+s}{U}\PY{l+s}{M}\PY{l+s}{P}\PY{l+s}{S}\PY{l+s}{ }\PY{l+s}{J}\PY{l+s}{O}\PY{l+s}{D}\PY{l+s}{S}\PY{l+s}{O}\PY{l+s}{U}\PY{l+s}{R}\PY{l+s}{C}\PY{l+s}{E}\PY{l+s}{ }\PY{l+s}{J}\PY{l+s}{O}\PY{l+s}{D}\PY{l+s}{T}\PY{l+s}{E}\PY{l+s}{S}\PY{l+s}{T}\PY{l+s}{\PYZsq{}}
        \PY{n+nv}{jodudefs} \PY{o}{=.} \PY{n+nv}{jpath}\PY{o}{\PYZam{}}\PY{o}{.}\PY{o}{\PYZgt{}} \PY{n+nv}{jodudirs}  \PY{c+c1}{NB. !(*)=. jpath}
        \PY{n+nv}{mask} \PY{o}{=.} \PY{n+nv}{jodudirs} \PY{o}{=} \PY{n+nv}{jodudefs}
        \PY{n+nl}{if.} \PY{l+m+mi}{1} \PY{n+nv}{e}\PY{o}{.} \PY{n+nv}{mask} \PY{n+nl}{do.}
          \PY{o}{\PYZgt{}} \PY{l+s}{\PYZsq{}}\PY{l+s}{m}\PY{l+s}{i}\PY{l+s}{s}\PY{l+s}{s}\PY{l+s}{i}\PY{l+s}{n}\PY{l+s}{g}\PY{l+s}{ }\PY{l+s}{J}\PY{l+s}{O}\PY{l+s}{D}\PY{l+s}{ }\PY{l+s}{f}\PY{l+s}{o}\PY{l+s}{l}\PY{l+s}{d}\PY{l+s}{e}\PY{l+s}{r}\PY{l+s}{s}\PY{l+s}{ }\PY{l+s}{\PYZhy{}}\PY{l+s}{ }\PY{l+s}{d}\PY{l+s}{e}\PY{l+s}{f}\PY{l+s}{i}\PY{l+s}{n}\PY{l+s}{e}\PY{l+s}{ }\PY{l+s}{b}\PY{l+s}{e}\PY{l+s}{f}\PY{l+s}{o}\PY{l+s}{r}\PY{l+s}{e}\PY{l+s}{ }\PY{l+s}{r}\PY{l+s}{u}\PY{l+s}{n}\PY{l+s}{n}\PY{l+s}{i}\PY{l+s}{n}\PY{l+s}{g}\PY{l+s}{ }\PY{l+s}{J}\PY{l+s}{O}\PY{l+s}{D}\PY{l+s}{ }\PY{l+s}{l}\PY{l+s}{a}\PY{l+s}{b}\PY{l+s}{s}\PY{l+s}{\PYZsq{}} \PY{o}{;} \PY{n+nv}{mask} \PY{o}{\PYZsh{}} \PY{n+nv}{jodudirs}
        \PY{n+nl}{else.}
          \PY{o}{\PYZgt{}} \PY{l+s}{\PYZsq{}}\PY{l+s}{a}\PY{l+s}{l}\PY{l+s}{l}\PY{l+s}{ }\PY{l+s}{J}\PY{l+s}{O}\PY{l+s}{D}\PY{l+s}{ }\PY{l+s}{f}\PY{l+s}{o}\PY{l+s}{l}\PY{l+s}{d}\PY{l+s}{e}\PY{l+s}{r}\PY{l+s}{s}\PY{l+s}{ }\PY{l+s}{a}\PY{l+s}{r}\PY{l+s}{e}\PY{l+s}{ }\PY{l+s}{d}\PY{l+s}{e}\PY{l+s}{f}\PY{l+s}{i}\PY{l+s}{n}\PY{l+s}{e}\PY{l+s}{d}\PY{l+s}{\PYZsq{}} \PY{o}{;} \PY{n+nv}{jodudefs}
        \PY{n+nl}{end.}
        \PY{n+nl}{)}
        
        \PY{n+nv}{TestJODDirectories} \PY{l+m+mi}{0}
\end{Verbatim}


    \begin{Verbatim}[commandchars=\\\{\}]
all JOD folders are defined
c:/jod                     
c:/jod/joddumps            
c:/jodtest/labtesting      
c:/jodtest/test            

    \end{Verbatim}

    \subsubsection{Lab cleanup}\label{lab-cleanup}

\textbf{WARNING: DICTIONARIES CREATED BY PRIOR RUNS OF LAB THIS WILL BE
DELETED IN THE NEXT STEP.}

When the Best Practices lab runs it creates a number of dictionaries in
the (\texttt{\textasciitilde{}JODSOURCE}) directory. The next step will
remove any of these dictionaries. \textbf{IF YOU CARE ABOUT THESE
DICTIONARIES STOP NOW.}

    \begin{Verbatim}[commandchars=\\\{\}]
{\color{incolor}In [{\color{incolor}4}]:} \PY{n+nv}{RemoveLabBestDictionaries\PYZus{}ijod\PYZus{}}\PY{o}{=:} \PY{n+nf}{3 : 0}
        \PY{n+nl}{if.} \PY{n+nv}{IFWIN} \PY{n+nl}{do.}
          \PY{n+nv}{shell} \PY{l+s}{\PYZsq{}}\PY{l+s}{r}\PY{l+s}{d}\PY{l+s}{ }\PY{l+s}{/}\PY{l+s}{s}\PY{l+s}{ }\PY{l+s}{/}\PY{l+s}{q}\PY{l+s}{ }\PY{l+s}{\PYZdq{}}\PY{l+s}{\PYZsq{}}\PY{o}{,}\PY{p}{(}\PY{n+nv}{jpath} \PY{l+s}{\PYZsq{}}\PY{l+s}{\PYZti{}}\PY{l+s}{J}\PY{l+s}{O}\PY{l+s}{D}\PY{l+s}{S}\PY{l+s}{O}\PY{l+s}{U}\PY{l+s}{R}\PY{l+s}{C}\PY{l+s}{E}\PY{l+s}{\PYZsq{}}\PY{p}{)}\PY{o}{,}\PY{l+s}{\PYZsq{}}\PY{l+s}{\PYZdq{}}\PY{l+s}{\PYZsq{}}
          \PY{n+nv}{smoutput} \PY{l+s}{\PYZsq{}}\PY{l+s}{L}\PY{l+s}{a}\PY{l+s}{b}\PY{l+s}{ }\PY{l+s}{t}\PY{l+s}{e}\PY{l+s}{m}\PY{l+s}{p}\PY{l+s}{o}\PY{l+s}{r}\PY{l+s}{a}\PY{l+s}{r}\PY{l+s}{y}\PY{l+s}{ }\PY{l+s}{b}\PY{l+s}{e}\PY{l+s}{s}\PY{l+s}{t}\PY{l+s}{ }\PY{l+s}{p}\PY{l+s}{r}\PY{l+s}{a}\PY{l+s}{c}\PY{l+s}{t}\PY{l+s}{i}\PY{l+s}{c}\PY{l+s}{e}\PY{l+s}{s}\PY{l+s}{ }\PY{l+s}{(}\PY{l+s}{w}\PY{l+s}{i}\PY{l+s}{n}\PY{l+s}{)}\PY{l+s}{ }\PY{l+s}{d}\PY{l+s}{i}\PY{l+s}{c}\PY{l+s}{t}\PY{l+s}{i}\PY{l+s}{o}\PY{l+s}{n}\PY{l+s}{a}\PY{l+s}{r}\PY{l+s}{i}\PY{l+s}{e}\PY{l+s}{s}\PY{l+s}{ }\PY{l+s}{e}\PY{l+s}{r}\PY{l+s}{a}\PY{l+s}{s}\PY{l+s}{e}\PY{l+s}{d}\PY{l+s}{\PYZsq{}}
        \PY{n+nl}{elseif.} \PY{n+nv}{IFUNIX} \PY{n+nl}{do.}
          \PY{c+c1}{NB. avoid blanks in Mac and Linux paths}
          \PY{n+nv}{shell} \PY{l+s}{\PYZsq{}}\PY{l+s}{r}\PY{l+s}{m}\PY{l+s}{ }\PY{l+s}{\PYZhy{}}\PY{l+s}{r}\PY{l+s}{f}\PY{l+s}{ }\PY{l+s}{\PYZsq{}}\PY{o}{,}\PY{n+nv}{jpath} \PY{l+s}{\PYZsq{}}\PY{l+s}{\PYZti{}}\PY{l+s}{J}\PY{l+s}{O}\PY{l+s}{D}\PY{l+s}{S}\PY{l+s}{O}\PY{l+s}{U}\PY{l+s}{R}\PY{l+s}{C}\PY{l+s}{E}\PY{l+s}{\PYZsq{}}
          \PY{n+nv}{smoutput} \PY{l+s}{\PYZsq{}}\PY{l+s}{L}\PY{l+s}{a}\PY{l+s}{b}\PY{l+s}{ }\PY{l+s}{t}\PY{l+s}{e}\PY{l+s}{m}\PY{l+s}{p}\PY{l+s}{o}\PY{l+s}{r}\PY{l+s}{a}\PY{l+s}{r}\PY{l+s}{y}\PY{l+s}{ }\PY{l+s}{b}\PY{l+s}{e}\PY{l+s}{s}\PY{l+s}{t}\PY{l+s}{ }\PY{l+s}{p}\PY{l+s}{r}\PY{l+s}{a}\PY{l+s}{c}\PY{l+s}{t}\PY{l+s}{i}\PY{l+s}{c}\PY{l+s}{e}\PY{l+s}{s}\PY{l+s}{ }\PY{l+s}{(}\PY{l+s}{m}\PY{l+s}{a}\PY{l+s}{c}\PY{l+s}{/}\PY{l+s}{l}\PY{l+s}{i}\PY{l+s}{n}\PY{l+s}{u}\PY{l+s}{x}\PY{l+s}{)}\PY{l+s}{ }\PY{l+s}{d}\PY{l+s}{i}\PY{l+s}{c}\PY{l+s}{t}\PY{l+s}{i}\PY{l+s}{o}\PY{l+s}{n}\PY{l+s}{a}\PY{l+s}{r}\PY{l+s}{i}\PY{l+s}{e}\PY{l+s}{s}\PY{l+s}{ }\PY{l+s}{e}\PY{l+s}{r}\PY{l+s}{a}\PY{l+s}{s}\PY{l+s}{e}\PY{l+s}{d}\PY{l+s}{\PYZsq{}}
        \PY{n+nl}{elseif.}\PY{n+nl}{do.}
          \PY{n+nv}{smoutput} \PY{l+s}{\PYZsq{}}\PY{l+s}{E}\PY{l+s}{r}\PY{l+s}{a}\PY{l+s}{s}\PY{l+s}{e}\PY{l+s}{ }\PY{l+s}{a}\PY{l+s}{n}\PY{l+s}{y}\PY{l+s}{ }\PY{l+s}{p}\PY{l+s}{r}\PY{l+s}{e}\PY{l+s}{v}\PY{l+s}{i}\PY{l+s}{o}\PY{l+s}{u}\PY{l+s}{s}\PY{l+s}{ }\PY{l+s}{t}\PY{l+s}{e}\PY{l+s}{m}\PY{l+s}{p}\PY{l+s}{o}\PY{l+s}{r}\PY{l+s}{a}\PY{l+s}{r}\PY{l+s}{y}\PY{l+s}{ }\PY{l+s}{b}\PY{l+s}{e}\PY{l+s}{s}\PY{l+s}{t}\PY{l+s}{ }\PY{l+s}{p}\PY{l+s}{r}\PY{l+s}{a}\PY{l+s}{c}\PY{l+s}{t}\PY{l+s}{i}\PY{l+s}{c}\PY{l+s}{e}\PY{l+s}{s}\PY{l+s}{ }\PY{l+s}{d}\PY{l+s}{i}\PY{l+s}{c}\PY{l+s}{t}\PY{l+s}{i}\PY{l+s}{o}\PY{l+s}{n}\PY{l+s}{a}\PY{l+s}{r}\PY{l+s}{i}\PY{l+s}{e}\PY{l+s}{s}\PY{l+s}{ }\PY{l+s}{m}\PY{l+s}{a}\PY{l+s}{n}\PY{l+s}{u}\PY{l+s}{a}\PY{l+s}{l}\PY{l+s}{l}\PY{l+s}{y}\PY{l+s}{.}\PY{l+s}{\PYZsq{}}
        \PY{n+nl}{end.}
        \PY{n+nl}{)}
\end{Verbatim}


    Remove Best Practices lab dictionaries.

    \begin{Verbatim}[commandchars=\\\{\}]
{\color{incolor}In [{\color{incolor}5}]:} \PY{c+c1}{NB. clear any previously opened dictionaries in master file}
        \PY{n+nv}{dpset} \PY{l+s}{\PYZsq{}}\PY{l+s}{R}\PY{l+s}{E}\PY{l+s}{S}\PY{l+s}{E}\PY{l+s}{T}\PY{l+s}{M}\PY{l+s}{E}\PY{l+s}{\PYZsq{}}
        
        \PY{c+c1}{NB. unregister dictionaries \PYZhy{} IGNORE ERRORS}
        \PY{l+m+mi}{3} \PY{n+nv}{regd}\PY{o}{\PYZam{}}\PY{o}{\PYZgt{}} \PY{o}{;}\PY{o}{:}\PY{l+s}{\PYZsq{}}\PY{l+s}{b}\PY{l+s}{p}\PY{l+s}{c}\PY{l+s}{o}\PY{l+s}{p}\PY{l+s}{y}\PY{l+s}{ }\PY{l+s}{b}\PY{l+s}{p}\PY{l+s}{t}\PY{l+s}{e}\PY{l+s}{s}\PY{l+s}{t}\PY{l+s}{\PYZsq{}} \PY{o}{[} \PY{l+m+mi}{3} \PY{n+nv}{od} \PY{l+s}{\PYZsq{}}\PY{l+s}{\PYZsq{}}
        
        \PY{c+c1}{NB. delete dictionary files}
        \PY{n+nv}{RemoveLabBestDictionaries} \PY{l+m+mi}{0}
\end{Verbatim}


    \begin{Verbatim}[commandchars=\\\{\}]
Lab temporary best practices (win) dictionaries erased

    \end{Verbatim}

    \subsubsection{JOD does not belong in the J
tree}\label{jod-does-not-belong-in-the-j-tree}

My first and most important bit of advice is simply:

\textbf{NEVER NEVER NEVER store your JOD dictionaries in J install
directories!}

I would also avoid any OS managed directories like

\begin{verbatim}
c:\Progam Files\..
\end{verbatim}

Create a JOD master dictionary directory root that is completely
independent of J. It's also a good idea to define a subdirectory
structure that mirrors J's versions. JOD creates binary \texttt{jfiles}.
These files are fairly stable but binaries can change when J changes.

    \begin{Verbatim}[commandchars=\\\{\}]
{\color{incolor}In [{\color{incolor}6}]:} \PY{c+c1}{NB. create a master JOD directory root outside of J\PYZsq{}s directories.}
        \PY{c+c1}{NB. This creation depends on a configured directory (\PYZti{}JODSOURCE).}
        \PY{n+nv}{smoutput} \PY{n+nv}{newd} \PY{l+s}{\PYZsq{}}\PY{l+s}{b}\PY{l+s}{p}\PY{l+s}{t}\PY{l+s}{e}\PY{l+s}{s}\PY{l+s}{t}\PY{l+s}{\PYZsq{}}\PY{o}{;}\PY{p}{(}\PY{n+nv}{jpath} \PY{l+s}{\PYZsq{}}\PY{l+s}{\PYZti{}}\PY{l+s}{J}\PY{l+s}{O}\PY{l+s}{D}\PY{l+s}{S}\PY{l+s}{O}\PY{l+s}{U}\PY{l+s}{R}\PY{l+s}{C}\PY{l+s}{E}\PY{l+s}{/}\PY{l+s}{b}\PY{l+s}{p}\PY{l+s}{t}\PY{l+s}{e}\PY{l+s}{s}\PY{l+s}{t}\PY{l+s}{\PYZsq{}}\PY{p}{)}\PY{o}{;}\PY{l+s}{\PYZsq{}}\PY{l+s}{b}\PY{l+s}{e}\PY{l+s}{s}\PY{l+s}{t}\PY{l+s}{ }\PY{l+s}{p}\PY{l+s}{r}\PY{l+s}{a}\PY{l+s}{c}\PY{l+s}{t}\PY{l+s}{i}\PY{l+s}{c}\PY{l+s}{e}\PY{l+s}{s}\PY{l+s}{ }\PY{l+s}{t}\PY{l+s}{e}\PY{l+s}{s}\PY{l+s}{t}\PY{l+s}{ }\PY{l+s}{d}\PY{l+s}{i}\PY{l+s}{c}\PY{l+s}{t}\PY{l+s}{i}\PY{l+s}{o}\PY{l+s}{n}\PY{l+s}{a}\PY{l+s}{r}\PY{l+s}{y}\PY{l+s}{\PYZsq{}}
        
        \PY{c+c1}{NB. This is an example \PYZhy{} use another root for your dictionaries.}
\end{Verbatim}


    \begin{Verbatim}[commandchars=\\\{\}]
+-+---------------------+------+-----------------------------+
|1|dictionary created ->|bptest|c:/jodtest/labtesting/bptest/|
+-+---------------------+------+-----------------------------+

    \end{Verbatim}

    \subsubsection{Load lab dictionary from dump
script}\label{load-lab-dictionary-from-dump-script}

To illustrate key features of JOD we need a nonempty dictionary. This
next step loads the (\texttt{bptest}) dictionary from a dump script
distributed with JOD.

    \begin{Verbatim}[commandchars=\\\{\}]
{\color{incolor}In [{\color{incolor}7}]:} \PY{c+c1}{NB. open new dictionary}
        \PY{n+nv}{od} \PY{l+s}{\PYZsq{}}\PY{l+s}{b}\PY{l+s}{p}\PY{l+s}{t}\PY{l+s}{e}\PY{l+s}{s}\PY{l+s}{t}\PY{l+s}{\PYZsq{}} \PY{o}{[} \PY{l+m+mi}{3} \PY{n+nv}{od} \PY{l+s}{\PYZsq{}}\PY{l+s}{\PYZsq{}}
        
        \PY{c+c1}{NB. load from example dump script}
        \PY{l+m+mi}{0}\PY{o}{!}\PY{o}{:}\PY{l+m+mi}{0} \PY{o}{\PYZlt{}}\PY{n+nv}{jpath} \PY{l+s}{\PYZsq{}}\PY{l+s}{\PYZti{}}\PY{l+s}{a}\PY{l+s}{d}\PY{l+s}{d}\PY{l+s}{o}\PY{l+s}{n}\PY{l+s}{s}\PY{l+s}{/}\PY{l+s}{g}\PY{l+s}{e}\PY{l+s}{n}\PY{l+s}{e}\PY{l+s}{r}\PY{l+s}{a}\PY{l+s}{l}\PY{l+s}{/}\PY{l+s}{j}\PY{l+s}{o}\PY{l+s}{d}\PY{l+s}{/}\PY{l+s}{j}\PY{l+s}{o}\PY{l+s}{d}\PY{l+s}{l}\PY{l+s}{a}\PY{l+s}{b}\PY{l+s}{s}\PY{l+s}{/}\PY{l+s}{l}\PY{l+s}{a}\PY{l+s}{b}\PY{l+s}{d}\PY{l+s}{u}\PY{l+s}{m}\PY{l+s}{p}\PY{l+s}{.}\PY{l+s}{i}\PY{l+s}{j}\PY{l+s}{s}\PY{l+s}{\PYZsq{}}
        
        \PY{c+c1}{NB. regenerate all references \PYZhy{} show last three messages}
        \PY{l+m+mi}{\PYZus{}3} \PY{o}{\PYZob{}}\PY{o}{.} \PY{l+m+mi}{0} \PY{n+nv}{globs}\PY{o}{\PYZam{}}\PY{o}{\PYZgt{}} \PY{o}{\PYZcb{}}\PY{o}{.} \PY{n+nv}{revo}\PY{l+s}{\PYZsq{}}\PY{l+s}{\PYZsq{}}
\end{Verbatim}


    \begin{Verbatim}[commandchars=\\\{\}]
+-+-------------------+------+
|1|1 word(s) put in ->|bptest|
+-+-------------------+------+
+-+--------------------+------+
|1|35 word(s) put in ->|bptest|
+-+--------------------+------+
+-+--------------------------------+------+
|1|36 word explanation(s) put in ->|bptest|
+-+--------------------------------+------+
+-+----------------------------+------+
|1|2 word document(s) put in ->|bptest|
+-+----------------------------+------+
+-+-------------------------+------+
|1|group <bstats> put in -> |bptest|
+-+-------------------------+------+
|1|group <sunmoon> put in ->|bptest|
+-+-------------------------+------+
NB. end-of-JOD-dump-file regenerate cross references with:  0 globs\&> \}. revo '' 
+-+-------------------------------------+------+
|1|<var> references put in ->           |bptest|
+-+-------------------------------------+------+
|1|<yeardates> references put in ->     |bptest|
+-+-------------------------------------+------+
|1|<NORISESET> is a noun - no references|      |
+-+-------------------------------------+------+

    \end{Verbatim}

    \subsubsection{Backup backup backup}\label{backup-backup-backup}

A wise man once said, \emph{"You're either backed up or f\&\%ked up!"}

\emph{Care to go over the options again?}

It is literally a snap to make a backup with JOD. So backup often!

    \begin{Verbatim}[commandchars=\\\{\}]
{\color{incolor}In [{\color{incolor}8}]:} \PY{c+c1}{NB. open the best practice dictionary}
        \PY{n+nv}{od} \PY{l+s}{\PYZsq{}}\PY{l+s}{b}\PY{l+s}{p}\PY{l+s}{t}\PY{l+s}{e}\PY{l+s}{s}\PY{l+s}{t}\PY{l+s}{\PYZsq{}} \PY{o}{[} \PY{l+m+mi}{3} \PY{n+nv}{od} \PY{l+s}{\PYZsq{}}\PY{l+s}{\PYZsq{}}
        
        \PY{c+c1}{NB. back it up}
        \PY{n+nv}{smoutput} \PY{n+nv}{packd} \PY{l+s}{\PYZsq{}}\PY{l+s}{b}\PY{l+s}{p}\PY{l+s}{t}\PY{l+s}{e}\PY{l+s}{s}\PY{l+s}{t}\PY{l+s}{\PYZsq{}}
        
        \PY{c+c1}{NB. restd recovers the most current backup}
        \PY{n+nv}{restd} \PY{l+s}{\PYZsq{}}\PY{l+s}{b}\PY{l+s}{p}\PY{l+s}{t}\PY{l+s}{e}\PY{l+s}{s}\PY{l+s}{t}\PY{l+s}{\PYZsq{}}
\end{Verbatim}


    \begin{Verbatim}[commandchars=\\\{\}]
+-+--------------------+------+-+
|1|dictionary packed ->|bptest|0|
+-+--------------------+------+-+
+-+----------------------+------+-+
|1|dictionary restored ->|bptest|0|
+-+----------------------+------+-+

    \end{Verbatim}

    \subsubsection{Take a script dump}\label{take-a-script-dump}

(\texttt{packd}) backs up binary \texttt{jfiles} but it's also a good
idea to "dump" your dictionaries as plain text. JOD can dump all open
dictionaries as a single J script. Script dumps are the most stable way
to store J dictionaries. The \texttt{jodsource} addon distributes all
JOD source code in this form.

    \begin{Verbatim}[commandchars=\\\{\}]
{\color{incolor}In [{\color{incolor}9}]:} \PY{c+c1}{NB. dump only (bptest)}
        \PY{n+nv}{od} \PY{l+s}{\PYZsq{}}\PY{l+s}{b}\PY{l+s}{p}\PY{l+s}{t}\PY{l+s}{e}\PY{l+s}{s}\PY{l+s}{t}\PY{l+s}{\PYZsq{}} \PY{o}{[} \PY{l+m+mi}{3} \PY{n+nv}{od} \PY{l+s}{\PYZsq{}}\PY{l+s}{\PYZsq{}}
        
        \PY{c+c1}{NB. (make) creates a dictionary dump in the dump subdirectory}
        \PY{n+nv}{bptestdump}\PY{o}{=:} \PY{n+nv}{showpass\PYZus{}ajod\PYZus{}} \PY{n+nv}{make} \PY{l+s}{\PYZsq{}}\PY{l+s}{\PYZsq{}}
\end{Verbatim}


    \begin{Verbatim}[commandchars=\\\{\}]
+-+---------------------------+--------------------------------------------+
|1|object(s) on path dumped ->|c:/jodtest/labtesting/bptest/dump/bptest.ijs|
+-+---------------------------+--------------------------------------------+

    \end{Verbatim}

    \subsubsection{Some uses of dump
scripts}\label{some-uses-of-dump-scripts}

JOD dump scripts can be used to reload, copy and merge dictionaries.

    \begin{Verbatim}[commandchars=\\\{\}]
{\color{incolor}In [{\color{incolor}10}]:} \PY{c+c1}{NB. reload bptest}
         \PY{n+nv}{od} \PY{l+s}{\PYZsq{}}\PY{l+s}{b}\PY{l+s}{p}\PY{l+s}{t}\PY{l+s}{e}\PY{l+s}{s}\PY{l+s}{t}\PY{l+s}{\PYZsq{}} \PY{o}{[} \PY{l+m+mi}{3} \PY{n+nv}{od} \PY{l+s}{\PYZsq{}}\PY{l+s}{\PYZsq{}}
         \PY{l+m+mi}{0}\PY{o}{!}\PY{o}{:}\PY{l+m+mi}{0} \PY{o}{\PYZob{}}\PY{o}{:}\PY{n+nv}{bptestdump}
\end{Verbatim}


    \begin{Verbatim}[commandchars=\\\{\}]
+-+-------------------+------+
|1|1 word(s) put in ->|bptest|
+-+-------------------+------+
+-+--------------------+------+
|1|35 word(s) put in ->|bptest|
+-+--------------------+------+
+-+--------------------------------+------+
|1|36 word explanation(s) put in ->|bptest|
+-+--------------------------------+------+
+-+----------------------------+------+
|1|2 word document(s) put in ->|bptest|
+-+----------------------------+------+
+-+-------------------------+------+
|1|group <bstats> put in -> |bptest|
+-+-------------------------+------+
|1|group <sunmoon> put in ->|bptest|
+-+-------------------------+------+
+-+------------------------------+------+
|1|dictionary document updated ->|bptest|
+-+------------------------------+------+
NB. end-of-JOD-dump-file regenerate cross references with:  0 globs\&> \}. revo '' 

    \end{Verbatim}

    Copy a dictionary.

    \begin{Verbatim}[commandchars=\\\{\}]
{\color{incolor}In [{\color{incolor}11}]:} \PY{c+c1}{NB. copy/merge bptest dictionary}
         \PY{n+nv}{smoutput} \PY{n+nv}{newd} \PY{l+s}{\PYZsq{}}\PY{l+s}{b}\PY{l+s}{p}\PY{l+s}{c}\PY{l+s}{o}\PY{l+s}{p}\PY{l+s}{y}\PY{l+s}{\PYZsq{}}\PY{o}{;}\PY{n+nv}{jpath} \PY{l+s}{\PYZsq{}}\PY{l+s}{\PYZti{}}\PY{l+s}{J}\PY{l+s}{O}\PY{l+s}{D}\PY{l+s}{S}\PY{l+s}{O}\PY{l+s}{U}\PY{l+s}{R}\PY{l+s}{C}\PY{l+s}{E}\PY{l+s}{/}\PY{l+s}{b}\PY{l+s}{p}\PY{l+s}{c}\PY{l+s}{o}\PY{l+s}{p}\PY{l+s}{y}\PY{l+s}{\PYZsq{}}
         \PY{n+nv}{od} \PY{l+s}{\PYZsq{}}\PY{l+s}{b}\PY{l+s}{p}\PY{l+s}{c}\PY{l+s}{o}\PY{l+s}{p}\PY{l+s}{y}\PY{l+s}{\PYZsq{}} \PY{o}{[} \PY{l+m+mi}{3} \PY{n+nv}{od} \PY{l+s}{\PYZsq{}}\PY{l+s}{\PYZsq{}}
         \PY{l+m+mi}{0}\PY{o}{!}\PY{o}{:}\PY{l+m+mi}{0} \PY{o}{\PYZob{}}\PY{o}{:}\PY{n+nv}{bptestdump}
         
         
         \PY{c+c1}{NB. clear path}
         \PY{n+nv}{dpset} \PY{l+s}{\PYZsq{}}\PY{l+s}{C}\PY{l+s}{L}\PY{l+s}{E}\PY{l+s}{A}\PY{l+s}{R}\PY{l+s}{P}\PY{l+s}{A}\PY{l+s}{T}\PY{l+s}{H}\PY{l+s}{\PYZsq{}}
\end{Verbatim}


    \begin{Verbatim}[commandchars=\\\{\}]
+-+---------------------+------+-----------------------------+
|1|dictionary created ->|bpcopy|c:/jodtest/labtesting/bpcopy/|
+-+---------------------+------+-----------------------------+
+-+-------------------+------+
|1|1 word(s) put in ->|bpcopy|
+-+-------------------+------+
+-+--------------------+------+
|1|35 word(s) put in ->|bpcopy|
+-+--------------------+------+
+-+--------------------------------+------+
|1|36 word explanation(s) put in ->|bpcopy|
+-+--------------------------------+------+
+-+----------------------------+------+
|1|2 word document(s) put in ->|bpcopy|
+-+----------------------------+------+
+-+-------------------------+------+
|1|group <bstats> put in -> |bpcopy|
+-+-------------------------+------+
|1|group <sunmoon> put in ->|bpcopy|
+-+-------------------------+------+
+-+------------------------------+------+
|1|dictionary document updated ->|bpcopy|
+-+------------------------------+------+
NB. end-of-JOD-dump-file regenerate cross references with:  0 globs\&> \}. revo '' 
+-+---------------+------+
|1|path cleared ->|bpcopy|
+-+---------------+------+

    \end{Verbatim}

    \subsubsection{Dump scripts are the best way to share and version
control
dictionaries}\label{dump-scripts-are-the-best-way-to-share-and-version-control-dictionaries}

Dump scripts can be used to share and version control dictionaries. See
this Git repository for examples.

\href{https://github.com/bakerjd99/joddumps}{Example JOD Dump Script
Repository}

    \subsubsection{Make a master re-register
script}\label{make-a-master-re-register-script}

JOD only sees the dictionaries registered in the \texttt{jmaster.ijf}
file so maintaining a list of registered dictionaries is a good idea.
JOD can generate a re-register script that you can edit.

Generate a re-register script and put it in your main JOD dictionary
directory root.

    \begin{Verbatim}[commandchars=\\\{\}]
{\color{incolor}In [{\color{incolor}12}]:} \PY{c+c1}{NB. generate re\PYZhy{}register script}
         \PY{n+nv}{rereg}\PY{o}{=:} \PY{n+nv}{showpass\PYZus{}ajod\PYZus{}} \PY{o}{;}\PY{o}{\PYZob{}}\PY{o}{:} \PY{l+m+mi}{5} \PY{n+nv}{od} \PY{l+s}{\PYZsq{}}\PY{l+s}{\PYZsq{}}
         
         \PY{c+c1}{NB. save it in the master root}
         \PY{n+nv}{rereg} \PY{n+nv}{write\PYZus{}ajod\PYZus{}} \PY{n+nv}{jpath} \PY{l+s}{\PYZsq{}}\PY{l+s}{\PYZti{}}\PY{l+s}{J}\PY{l+s}{O}\PY{l+s}{D}\PY{l+s}{S}\PY{l+s}{O}\PY{l+s}{U}\PY{l+s}{R}\PY{l+s}{C}\PY{l+s}{E}\PY{l+s}{/}\PY{l+s}{j}\PY{l+s}{o}\PY{l+s}{d}\PY{l+s}{r}\PY{l+s}{e}\PY{l+s}{g}\PY{l+s}{i}\PY{l+s}{s}\PY{l+s}{t}\PY{l+s}{e}\PY{l+s}{r}\PY{l+s}{.}\PY{l+s}{i}\PY{l+s}{j}\PY{l+s}{s}\PY{l+s}{\PYZsq{}}
\end{Verbatim}


    \begin{Verbatim}[commandchars=\\\{\}]
NB. JOD registrations: 09 Dec 2019 09:16:16
NB. require 'general/jod'
3 regd\&> \}. od'' [ 3 od ''
regd 'gps';'c:/jod/gps/'
regd 'image';'c:/jod/image/'
regd 'mwecc';'c:/jod/mwecc/'
regd 'play';'c:/jod/play/'
regd 'smug';'c:/jod/smug/'
regd 'smugdev';'c:/jod/smugdev/'
regd 'jodtest';'c:/jod/jodtest/'
regd 'utils';'c:/jod/utils/'
regd 'jod';'c:/jod/jod/'
regd 'joddev';'c:/jod/joddev/'
regd 'mep';'c:/jod/mep/'
regd 'fit';'c:/fit/jod/fit/'
regd 'fct';'c:/jod/fct/'
regd 'bitjd';'c:/jod/bitjd/'
regd 'bci';'c:/jod/bci/'
regd 'docs';'c:/jod/docs/'
regd 'imex';'c:/jod/imex/'
regd 'smugpyter';'c:/jod/smugpyter/'
regd 'bcidev';'c:/jod/bcidev/'
regd 'fd';'c:/jod/fd/'
regd 'testjod00';'c:/jodtest/test/testjod00/'
regd 'testjod01';'c:/jodtest/test/testjod01/'
regd 'testjod02';'c:/jodtest/test/testjod02/'
regd 'testjod03';'c:/jodtest/test/testjod03/'
regd 'testjod04';'c:/jodtest/test/testjod04/'
regd 'testjod05';'c:/jodtest/test/testjod05/'
regd 'testjod06';'c:/jodtest/test/testjod06/'
regd 'testjod07';'c:/jodtest/test/testjod07/'
regd 'testjod08';'c:/jodtest/test/testjod08/'
regd 'testjod09';'c:/jodtest/test/testjod09/'
regd 'testjod10';'c:/jodtest/test/testjod10/'
regd 'testjod11';'c:/jodtest/test/testjod11/'
regd 'testjod12';'c:/jodtest/test/testjod12/'
regd 'testjod13';'c:/jodtest/test/testjod13/'
regd 'testjod14';'c:/jodtest/test/testjod14/'
regd 'testjod15';'c:/jodtest/test/testjod15/'
regd 'testjod100';'c:/jodtest/test/testjod100/'
regd 'testjod1000';'c:/jodtest/test/testjod1000/'
regd 'rxsdemo200050854116529950128671240971375103355';'c:/users/john/j901-user/joddicts/rxsdemo200050854116529950128671240971375103355/'
regd 'rxsdemo24347018837875837655937079410780880785';'c:/users/john/j901-user/joddicts/rxsdemo24347018837875837655937079410780880785/'
regd 'utilsload';'c:/jodtest/labtesting/utilsload/'
regd 'jodload';'c:/jodtest/labtesting/jodload/'
regd 'joddevload';'c:/jodtest/labtesting/joddevload/'
regd 'lab';'c:/users/john/j901-user/joddicts/lab/'
regd 'labdev';'c:/users/john/j901-user/joddicts/labdev/'
regd 'toy';'c:/users/john/j901-user/joddicts/toy/'
regd 'playpen';'c:/users/john/j901-user/joddicts/playpen/'
regd 'bptest';'c:/jodtest/labtesting/bptest/'
regd 'bpcopy';'c:/jodtest/labtesting/bpcopy/'

    \end{Verbatim}

    \subsubsection{\texorpdfstring{Set library dictionaries to
\texttt{READONLY}}{Set library dictionaries to READONLY}}\label{set-library-dictionaries-to-readonly}

Open JOD dictionaries define a search path. The first dictionary on the
path is the only dictionary that can be changed. It is called the "put"
dictionary. Even though nonput dictionaries cannot be changed by JOD
it's a good idea to set them \texttt{READONLY} because:

\texttt{READONLY} dictionaries can be accessed by any number of JOD
tasks. \texttt{READWRITE} dictionaries can only be accessed by one task.

Keeping libraries \texttt{READONLY} prevents accidental put's as you
open and close dictionaries.

    \begin{Verbatim}[commandchars=\\\{\}]
{\color{incolor}In [{\color{incolor}13}]:} \PY{c+c1}{NB. make bptest READONLY}
         \PY{n+nv}{od} \PY{l+s}{\PYZsq{}}\PY{l+s}{b}\PY{l+s}{p}\PY{l+s}{t}\PY{l+s}{e}\PY{l+s}{s}\PY{l+s}{t}\PY{l+s}{\PYZsq{}} \PY{o}{[} \PY{l+m+mi}{3} \PY{n+nv}{od} \PY{l+s}{\PYZsq{}}\PY{l+s}{\PYZsq{}}
         \PY{n+nv}{dpset} \PY{l+s}{\PYZsq{}}\PY{l+s}{R}\PY{l+s}{E}\PY{l+s}{A}\PY{l+s}{D}\PY{l+s}{O}\PY{l+s}{N}\PY{l+s}{L}\PY{l+s}{Y}\PY{l+s}{\PYZsq{}}
         
         \PY{c+c1}{NB. bpcopy is now the put dictionary}
         \PY{n+nv}{od} \PY{o}{;}\PY{o}{:}\PY{l+s}{\PYZsq{}}\PY{l+s}{b}\PY{l+s}{p}\PY{l+s}{c}\PY{l+s}{o}\PY{l+s}{p}\PY{l+s}{y}\PY{l+s}{ }\PY{l+s}{b}\PY{l+s}{p}\PY{l+s}{t}\PY{l+s}{e}\PY{l+s}{s}\PY{l+s}{t}\PY{l+s}{\PYZsq{}} \PY{o}{[} \PY{l+m+mi}{3} \PY{n+nv}{od} \PY{l+s}{\PYZsq{}}\PY{l+s}{\PYZsq{}}
         
         \PY{c+c1}{NB. first group/suite sets path}
         \PY{n+nv}{grp} \PY{l+s}{\PYZsq{}}\PY{l+s}{a}\PY{l+s}{g}\PY{l+s}{r}\PY{l+s}{o}\PY{l+s}{u}\PY{l+s}{p}\PY{l+s}{\PYZsq{}}\PY{o}{;} \PY{o}{;}\PY{o}{:}\PY{l+s}{\PYZsq{}}\PY{l+s}{d}\PY{l+s}{a}\PY{l+s}{t}\PY{l+s}{e}\PY{l+s}{c}\PY{l+s}{h}\PY{l+s}{e}\PY{l+s}{c}\PY{l+s}{k}\PY{l+s}{ }\PY{l+s}{y}\PY{l+s}{e}\PY{l+s}{a}\PY{l+s}{r}\PY{l+s}{d}\PY{l+s}{a}\PY{l+s}{t}\PY{l+s}{e}\PY{l+s}{s}\PY{l+s}{ }\PY{l+s}{t}\PY{l+s}{o}\PY{l+s}{d}\PY{l+s}{a}\PY{l+s}{y}\PY{l+s}{ }\PY{l+s}{s}\PY{l+s}{u}\PY{l+s}{n}\PY{l+s}{r}\PY{l+s}{i}\PY{l+s}{s}\PY{l+s}{e}\PY{l+s}{s}\PY{l+s}{e}\PY{l+s}{t}\PY{l+s}{0}\PY{l+s}{\PYZsq{}}
         
         \PY{c+c1}{NB. note dictionary path}
         \PY{n+nv}{did} \PY{o}{\PYZti{}} \PY{l+m+mi}{0}
\end{Verbatim}


    \begin{Verbatim}[commandchars=\\\{\}]
+-+-------------------------------------------------------------+
|1|+------+--+-----+-----+-------+-------+------+--------------+|
| ||      |--|Words|Tests|Groups*|Suites*|Macros|Path*         ||
| |+------+--+-----+-----+-------+-------+------+--------------+|
| ||bpcopy|rw|36   |0    |3      |0      |0     |/bpcopy/bptest||
| |+------+--+-----+-----+-------+-------+------+--------------+|
| ||bptest|ro|36   |0    |2      |0      |0     |/bptest       ||
| |+------+--+-----+-----+-------+-------+------+--------------+|
+-+-------------------------------------------------------------+

    \end{Verbatim}

    \subsubsection{Keep references updated}\label{keep-references-updated}

JOD stores word references. References enable many useful operations.
References allow (\texttt{getrx}) to load words that call other words in
new contexts.

    \begin{Verbatim}[commandchars=\\\{\}]
{\color{incolor}In [{\color{incolor}14}]:} \PY{c+c1}{NB. only put dictionary references need updating \PYZhy{} show last five messages}
         \PY{l+m+mi}{\PYZus{}5} \PY{o}{\PYZob{}}\PY{o}{.} \PY{l+m+mi}{0} \PY{n+nv}{globs}\PY{o}{\PYZam{}}\PY{o}{\PYZgt{}} \PY{o}{\PYZcb{}}\PY{o}{.} \PY{n+nv}{revo} \PY{l+s}{\PYZsq{}}\PY{l+s}{\PYZsq{}}
\end{Verbatim}


    \begin{Verbatim}[commandchars=\\\{\}]
+-+-------------------------------------+------+
|1|<tan> references put in ->           |bpcopy|
+-+-------------------------------------+------+
|1|<today> references put in ->         |bpcopy|
+-+-------------------------------------+------+
|1|<var> references put in ->           |bpcopy|
+-+-------------------------------------+------+
|1|<yeardates> references put in ->     |bpcopy|
+-+-------------------------------------+------+
|1|<NORISESET> is a noun - no references|      |
+-+-------------------------------------+------+

    \end{Verbatim}

    \subsubsection{Document dictionary
objects}\label{document-dictionary-objects}

Documentation is a long standing sore point for programmers. Most of
them hate it. Some claim it's unnecessary and even distracting. Others
put in half baked efforts. In my opinion this "isn't even wrong!" Good
documentation elevates code. In
\href{http://www.literateprogramming.com/knuthweb.pdf}{Knuth's opinion}
it separates "literate programming" from the alternative - "illiterate
programming."

JOD provides a number of easy ways to document code. You can enter a
single sentence or a large dissertation. I would recommend the former.
See JOD documentation for more documentation options.

    \begin{Verbatim}[commandchars=\\\{\}]
{\color{incolor}In [{\color{incolor}15}]:} \PY{c+c1}{NB. for new words try a single line of documentation.}
         \PY{n+nv}{afterlaststr}\PY{o}{=:}\PY{o}{]} \PY{o}{\PYZcb{}}\PY{o}{.}\PY{o}{\PYZti{}} \PY{o}{\PYZsh{}}\PY{o}{@}\PY{o}{[} \PY{o}{+} \PY{l+m+mi}{1}\PY{o}{\PYZam{}}\PY{p}{(}\PY{n+nv}{i}\PY{o}{:}\PY{o}{\PYZti{}}\PY{p}{)}\PY{o}{@}\PY{p}{(}\PY{o}{[} \PY{n+nv}{E}\PY{o}{.} \PY{o}{]}\PY{p}{)}
         \PY{n+nv}{put} \PY{l+s}{\PYZsq{}}\PY{l+s}{a}\PY{l+s}{f}\PY{l+s}{t}\PY{l+s}{e}\PY{l+s}{r}\PY{l+s}{l}\PY{l+s}{a}\PY{l+s}{s}\PY{l+s}{t}\PY{l+s}{s}\PY{l+s}{t}\PY{l+s}{r}\PY{l+s}{\PYZsq{}}
         
         \PY{c+c1}{NB. insert sentence}
         \PY{l+m+mi}{0} \PY{l+m+mi}{8} \PY{n+nv}{put} \PY{l+s}{\PYZsq{}}\PY{l+s}{a}\PY{l+s}{f}\PY{l+s}{t}\PY{l+s}{e}\PY{l+s}{r}\PY{l+s}{l}\PY{l+s}{a}\PY{l+s}{s}\PY{l+s}{t}\PY{l+s}{s}\PY{l+s}{t}\PY{l+s}{r}\PY{l+s}{\PYZsq{}}\PY{o}{;}\PY{l+s}{\PYZsq{}}\PY{l+s}{r}\PY{l+s}{e}\PY{l+s}{t}\PY{l+s}{a}\PY{l+s}{i}\PY{l+s}{n}\PY{l+s}{s}\PY{l+s}{ }\PY{l+s}{s}\PY{l+s}{t}\PY{l+s}{r}\PY{l+s}{i}\PY{l+s}{n}\PY{l+s}{g}\PY{l+s}{ }\PY{l+s}{(}\PY{l+s}{y}\PY{l+s}{)}\PY{l+s}{ }\PY{l+s}{a}\PY{l+s}{f}\PY{l+s}{t}\PY{l+s}{e}\PY{l+s}{r}\PY{l+s}{ }\PY{l+s}{l}\PY{l+s}{a}\PY{l+s}{s}\PY{l+s}{t}\PY{l+s}{ }\PY{l+s}{o}\PY{l+s}{c}\PY{l+s}{c}\PY{l+s}{u}\PY{l+s}{r}\PY{l+s}{r}\PY{l+s}{e}\PY{l+s}{n}\PY{l+s}{c}\PY{l+s}{e}\PY{l+s}{ }\PY{l+s}{o}\PY{l+s}{f}\PY{l+s}{ }\PY{l+s}{(}\PY{l+s}{x}\PY{l+s}{)}\PY{l+s}{\PYZsq{}}
\end{Verbatim}


    \begin{Verbatim}[commandchars=\\\{\}]
+-+-------------------------------+------+
|1|1 word explanation(s) put in ->|bpcopy|
+-+-------------------------------+------+

    \end{Verbatim}

    JOD uses documentation in new contexts.

    \begin{Verbatim}[commandchars=\\\{\}]
{\color{incolor}In [{\color{incolor}16}]:} \PY{c+c1}{NB. for tacits the docment sentence is fetched}
         \PY{n+nv}{smoutput} \PY{n+nv}{disp} \PY{l+s}{\PYZsq{}}\PY{l+s}{a}\PY{l+s}{f}\PY{l+s}{t}\PY{l+s}{e}\PY{l+s}{r}\PY{l+s}{l}\PY{l+s}{a}\PY{l+s}{s}\PY{l+s}{t}\PY{l+s}{s}\PY{l+s}{t}\PY{l+s}{r}\PY{l+s}{\PYZsq{}}
         
         \PY{c+c1}{NB. briefly describe sunmoon group  \PYZhy{} the payback for entering those sentences}
         \PY{n+nv}{hlpnl} \PY{o}{\PYZcb{}}\PY{o}{.} \PY{n+nv}{grp} \PY{l+s}{\PYZsq{}}\PY{l+s}{s}\PY{l+s}{u}\PY{l+s}{n}\PY{l+s}{m}\PY{l+s}{o}\PY{l+s}{o}\PY{l+s}{n}\PY{l+s}{\PYZsq{}}
\end{Verbatim}


    \begin{Verbatim}[commandchars=\\\{\}]
NB. retains string (y) after last occurrence of (x)
afterlaststr=:] \}.\textasciitilde{} \#@[ + 1\&(i:\textasciitilde{})@([ E. ])
+-----------+----------------------------------------------------------------------------+
|NORISESET  |indicates sun never rises or sets in (sunriseset0) and (sunriseset1) results|
|arctan     |arc tangent                                                                 |
|calmoons   |calendar dates of new and full moons                                        |
|cos        |cosine radians                                                              |
|datecheck  |checks dates in YYYY MM DD format                                           |
|fromjulian |converts Julian day numbers to dates, converse (tojulian)                   |
|moons      |times of new and full moons for n calendar years                            |
|round      |round y to nearest x (e.g. 1000 round 12345)                                |
|sin        |sine radians                                                                |
|sunriseset0|computes sun rise and set times - see long documentation                    |
|sunriseset1|computes sun rise and set times - see long documentation                    |
|tabit      |promotes only atoms and lists to tables                                     |
|tan        |tan radians                                                                 |
|today      |returns todays date                                                         |
|yeardates  |returns all valid dates for n calendar years                                |
+-----------+----------------------------------------------------------------------------+

    \end{Verbatim}

    An example of long documentation.

    \begin{Verbatim}[commandchars=\\\{\}]
{\color{incolor}In [{\color{incolor}17}]:} \PY{c+c1}{NB. long document}
         \PY{l+m+mi}{0} \PY{l+m+mi}{9} \PY{n+nv}{disp} \PY{l+s}{\PYZsq{}}\PY{l+s}{s}\PY{l+s}{u}\PY{l+s}{n}\PY{l+s}{r}\PY{l+s}{i}\PY{l+s}{s}\PY{l+s}{e}\PY{l+s}{s}\PY{l+s}{e}\PY{l+s}{t}\PY{l+s}{0}\PY{l+s}{\PYZsq{}}
\end{Verbatim}


    \begin{Verbatim}[commandchars=\\\{\}]
*sunriseset0 v-- sunrise and sunset times.

This  verb has been adapted from a BASIC program submitted by
Robin  G.  Stuart  Sky  \&  Telescope's  shortest  sunrise/set
program  contest. Winning  entries were listed  in the  March
1995 Astronomical Computing column.

The  J version of this algorithm has been vectorized.  It can
compute any number of sunrise and sunset times in one call.

NB. verbatim:

The (y) argument is a 5*n floating point table where:

 0\{ is latitude in degrees with northern latitudes positive.
 1\{ is longitude in degrees with western longitudes negative.
 2\{ is western time zones expressed as positive whole hours.
 3\{ is the month number.
 4\{ is the day number.

The result is a numeric table with four rows.  To handle the cases
when the sun never rises or sets the first two elements of the
corresponding result columns n are:

 0\{ is NORISESET an invalid hour indicating no rise or set
 1\{ is 0 when the sun never rises
 1\{ is 1 when the sun never sets

Warning: this algorithm breaks for latitudes close
to the South pole.

The original BASIC code has been slightly modified
to use control structures in place of GOTO's and line numbers.

Adapted from:

' Sunrise/set by R. G. Stuart,  Mexico City, Mexico
PI = 3.14159265\#: DR = PI / 180: RD = 1 / DR
INPUT "Lat, Long (deg)"; B5, L5
INPUT "Time zone (hrs)"; H
B5 = DR * B5
INPUT "Month, day"; M, D
N = INT(275 * M / 9) - 2 * INT((M + 9) / 12) + D - 30
L0 = 4.8771 + .0172 * (N + .5 - L5 / 360)
C = .03342 * SIN(L0 + 1.345)
C2 = RD * (ATN(TAN(L0 + C)) - ATN(.9175 * TAN(L0 + C)) - C)
SD = .3978 * SIN(L0 + C): CD = SQR(1 - SD * SD)
SC = (SD * SIN(B5) + .0145) / (COS(B5) * CD)
IF ABS(SC) <= 1 THEN
  C3 = RD * ATN(SC / SQR(1 - SC * SC))
  R1 = 6 - H - (L5 + C2 + C3) / 15
  HR = INT(R1): MR = INT((R1 - HR) * 60)
  PRINT USING "Sunrise at \#\#:\#\#"; HR; MR
  S1 = 18 - H - (L5 + C2 - C3) / 15
  HS = INT(S1): MS = INT((S1 - HS) * 60)
  PRINT USING "Sunset at \#\#:\#\#"; HS; MS
ELSEIF SC > 1 THEN
  PRINT "Sun up all day"
ELSEIF SC < -1 THEN
  PRINT "Sun down all day"
END IF
END


monad: ntRiseset =. sunriseset0 flBLHMD

  NB. rise and set times at Dog Lake today (daylight savings)
  td=. (44 + 19\%60),(- 76 + 21\%60), 4 , \}. today 0
  sunriseset0 td

  NB. rise and set times on June 30 on Greenwich meridian
  t0=.   0 0 0 6 30   NB. equator
  t1=.  49 0 0 6 30   NB. north - lat of western US/Canada border
  t2=. \_47 0 0 6 30   NB. south - southern Chile and Argentina
  t3=.  75 0 0 6 30   NB. far north (sun always up)
  t4=. \_75 0 0 6 30   NB. far south (sun always down)

  sunriseset0 t0

  sunriseset0 t0 , t1 , t2 , t3 ,: t4

  NB. times on equator for March 21 for all 1 hour time zones
  sunriseset0  0 0 ,"1 (,.i. 24) ,"1 ] 3 21

  NB. times for calendar year 1995 on the Greenwich meridian
  md95=.  47 0 0 ,"1 \}."1 yeardates 1995
  rs095=. sunriseset0 md95






    \end{Verbatim}

    \subsubsection{Define your own JOD
shortcuts}\label{define-your-own-jod-shortcuts}

JOD words can be used within arbitrary J programs. If you don't find a
JOD primitive that meets your needs do a little programming.

There are many examples of JOD programming in JOD's source code. Install
the \texttt{jodsource} addon to get JOD source code.

    \begin{Verbatim}[commandchars=\\\{\}]
{\color{incolor}In [{\color{incolor}18}]:} \PY{c+c1}{NB. examples of using JOD words to define new facilities}
         
         \PY{c+c1}{NB. describe a JOD group}
         \PY{n+nv}{hg\PYZus{}ijod\PYZus{}}\PY{o}{=:} \PY{o}{[}\PY{o}{:} \PY{n+nv}{hlpnl} \PY{o}{[}\PY{o}{:} \PY{o}{\PYZcb{}}\PY{o}{.} \PY{n+nv}{grp}
         
         \PY{c+c1}{NB. re\PYZhy{}reference put dictionary show any errors}
         \PY{n+nv}{reref\PYZus{}ijod\PYZus{}}\PY{o}{=:} \PY{l+m+mi}{3} \PY{o}{:} \PY{l+s}{\PYZsq{}}\PY{l+s}{(}\PY{l+s}{n}\PY{l+s}{,}\PY{l+s}{.}\PY{l+s}{s}\PY{l+s}{)}\PY{l+s}{ }\PY{l+s}{\PYZsh{}}\PY{l+s}{\PYZti{}}\PY{l+s}{ }\PY{l+s}{\PYZhy{}}\PY{l+s}{.}\PY{l+s}{;}\PY{l+s}{0}\PY{l+s}{\PYZob{}}\PY{l+s}{\PYZdq{}}\PY{l+s}{1}\PY{l+s}{ }\PY{l+s}{s}\PY{l+s}{=}\PY{l+s}{.}\PY{l+s}{0}\PY{l+s}{ }\PY{l+s}{g}\PY{l+s}{l}\PY{l+s}{o}\PY{l+s}{b}\PY{l+s}{s}\PY{l+s}{\PYZam{}}\PY{l+s}{\PYZgt{}}\PY{l+s}{n}\PY{l+s}{=}\PY{l+s}{.}\PY{l+s}{\PYZcb{}}\PY{l+s}{.}\PY{l+s}{r}\PY{l+s}{e}\PY{l+s}{v}\PY{l+s}{o}\PY{l+s}{\PYZsq{}\PYZsq{}}\PY{l+s}{\PYZsq{}\PYZsq{}}\PY{l+s}{ }\PY{l+s}{[}\PY{l+s}{ }\PY{l+s}{y}\PY{l+s}{\PYZsq{}}
         
         \PY{c+c1}{NB. show words referenced by words in a group that are not in the group}
         \PY{n+nv}{jodg\PYZus{}ijod\PYZus{}}\PY{o}{=:} \PY{l+s}{\PYZsq{}}\PY{l+s}{a}\PY{l+s}{g}\PY{l+s}{r}\PY{l+s}{o}\PY{l+s}{u}\PY{l+s}{p}\PY{l+s}{\PYZsq{}}
         \PY{n+nv}{nx\PYZus{}ijod\PYZus{}}\PY{o}{=:} \PY{l+m+mi}{3} \PY{o}{:} \PY{l+s}{\PYZsq{}}\PY{l+s}{(}\PY{l+s}{a}\PY{l+s}{l}\PY{l+s}{l}\PY{l+s}{r}\PY{l+s}{e}\PY{l+s}{f}\PY{l+s}{s}\PY{l+s}{ }\PY{l+s}{ }\PY{l+s}{\PYZcb{}}\PY{l+s}{.}\PY{l+s}{ }\PY{l+s}{g}\PY{l+s}{n}\PY{l+s}{)}\PY{l+s}{ }\PY{l+s}{\PYZhy{}}\PY{l+s}{.}\PY{l+s}{ }\PY{l+s}{g}\PY{l+s}{n}\PY{l+s}{=}\PY{l+s}{.}\PY{l+s}{ }\PY{l+s}{g}\PY{l+s}{r}\PY{l+s}{p}\PY{l+s}{ }\PY{l+s}{j}\PY{l+s}{o}\PY{l+s}{d}\PY{l+s}{g}\PY{l+s}{\PYZsq{}}
         
         \PY{c+c1}{NB. missing from (agroup)}
         \PY{n+nv}{nx} \PY{l+s}{\PYZsq{}}\PY{l+s}{\PYZsq{}}
\end{Verbatim}


    \begin{Verbatim}[commandchars=\\\{\}]
+---------+------+---+---+-----+---+
|NORISESET|arctan|cos|sin|tabit|tan|
+---------+------+---+---+-----+---+

    \end{Verbatim}

    \subsubsection{Customize JOD edit
facilities}\label{customize-jod-edit-facilities}

The main JOD edit words (\texttt{nw}) and (\texttt{ed}) can be
customized by defining a \texttt{DOCUMENTCOMMAND} script.

\emph{Note: Verbs that spawn J editors may not work in Jupyter labs. The
following (\texttt{nw}) call opens a JQT or JHS editor in standard J
front ends but does nothing here. This is because the J kernal is
essentially a barebones \texttt{jconsole.exe} process that is is not
running JQT.}

    \begin{Verbatim}[commandchars=\\\{\}]
{\color{incolor}In [{\color{incolor}19}]:} \PY{c+c1}{NB. define document command script \PYZhy{} \PYZob{}\PYZti{}N\PYZti{}\PYZcb{} is word name placeholder}
         \PY{n+nv}{DOCUMENTCOMMAND\PYZus{}ijod\PYZus{}}\PY{o}{=:} \PY{n+ni}{0 : 0}
         \PY{l+s}{s}\PY{l+s}{m}\PY{l+s}{o}\PY{l+s}{u}\PY{l+s}{t}\PY{l+s}{p}\PY{l+s}{u}\PY{l+s}{t}\PY{l+s}{ }\PY{l+s}{p}\PY{l+s}{r}\PY{l+s}{ }\PY{l+s}{\PYZsq{}}\PY{l+s}{\PYZob{}}\PY{l+s}{\PYZti{}}\PY{l+s}{N}\PY{l+s}{\PYZti{}}\PY{l+s}{\PYZcb{}}\PY{l+s}{\PYZsq{}}
         \PY{n+nl}{)}
         
         \PY{c+c1}{NB. create a new word \PYZhy{} opens an edit window in JQT and JHS }
         \PY{c+c1}{NB. nw \PYZsq{}bpword\PYZsq{}}
\end{Verbatim}


    Find the (\texttt{bpword}) edit window and note how
\texttt{DOCUMENTCOMMAND} text has been modified and insert. When the
script window is saved with CRTL-W \texttt{DOCUMENTCOMMAND} runs.

Run and close the (\texttt{bpword}) edit window.

Edit a word in the dictionary. \textbf{\emph{JOD is always editing
copies.}} Originals can only be changed with explicit (\texttt{put})
operations.

    \begin{Verbatim}[commandchars=\\\{\}]
{\color{incolor}In [{\color{incolor}20}]:} \PY{c+c1}{NB. load dictionary word into edit window \PYZhy{} requires JQT or JHS front ends}
         \PY{n+nv}{ed} \PY{l+s}{\PYZsq{}}\PY{l+s}{s}\PY{l+s}{u}\PY{l+s}{n}\PY{l+s}{r}\PY{l+s}{i}\PY{l+s}{s}\PY{l+s}{e}\PY{l+s}{s}\PY{l+s}{e}\PY{l+s}{t}\PY{l+s}{0}\PY{l+s}{\PYZsq{}}
         
         \PY{c+c1}{NB. find the (sunriseset0) edit window and note DOCUMENTCOMMAND }
         
         \PY{c+c1}{NB. close edit window}
\end{Verbatim}


    \begin{Verbatim}[commandchars=\\\{\}]
\$\$\$edit\$\$\$c:/users/john/j901-user/temp/sunriseset0.ijs

    \end{Verbatim}

    \subsubsection{Define JOD project
macros}\label{define-jod-project-macros}

When programming with JOD you typically open a set of dictionaries. Load
system scripts and define some handy nouns. This can be done in a
project macro script.

    \begin{Verbatim}[commandchars=\\\{\}]
{\color{incolor}In [{\color{incolor}21}]:} \PY{c+c1}{NB. define a project macro \PYZhy{} I use the prefix prj for such scripts}
         \PY{n+nv}{prjsunmoon}\PY{o}{=:} \PY{n+ni}{0 : 0}
         
         \PY{l+s}{N}\PY{l+s}{B}\PY{l+s}{.}\PY{l+s}{ }\PY{l+s}{s}\PY{l+s}{t}\PY{l+s}{a}\PY{l+s}{n}\PY{l+s}{d}\PY{l+s}{a}\PY{l+s}{r}\PY{l+s}{d}\PY{l+s}{ }\PY{l+s}{j}\PY{l+s}{ }\PY{l+s}{s}\PY{l+s}{c}\PY{l+s}{r}\PY{l+s}{i}\PY{l+s}{p}\PY{l+s}{t}\PY{l+s}{s}
         \PY{l+s}{r}\PY{l+s}{e}\PY{l+s}{q}\PY{l+s}{u}\PY{l+s}{i}\PY{l+s}{r}\PY{l+s}{e}\PY{l+s}{ }\PY{l+s}{\PYZsq{}}\PY{l+s}{d}\PY{l+s}{e}\PY{l+s}{b}\PY{l+s}{u}\PY{l+s}{g}\PY{l+s}{ }\PY{l+s}{t}\PY{l+s}{a}\PY{l+s}{s}\PY{l+s}{k}\PY{l+s}{\PYZsq{}}
         
         \PY{l+s}{N}\PY{l+s}{B}\PY{l+s}{.}\PY{l+s}{ }\PY{l+s}{l}\PY{l+s}{o}\PY{l+s}{c}\PY{l+s}{a}\PY{l+s}{l}\PY{l+s}{ }\PY{l+s}{s}\PY{l+s}{c}\PY{l+s}{r}\PY{l+s}{i}\PY{l+s}{p}\PY{l+s}{t}\PY{l+s}{ }\PY{l+s}{n}\PY{l+s}{o}\PY{l+s}{u}\PY{l+s}{n}\PY{l+s}{s}\PY{l+s}{ }
         \PY{l+s}{j}\PY{l+s}{o}\PY{l+s}{d}\PY{l+s}{g}\PY{l+s}{\PYZus{}}\PY{l+s}{i}\PY{l+s}{j}\PY{l+s}{o}\PY{l+s}{d}\PY{l+s}{\PYZus{}}\PY{l+s}{=}\PY{l+s}{:}\PY{l+s}{ }\PY{l+s}{\PYZsq{}}\PY{l+s}{s}\PY{l+s}{u}\PY{l+s}{n}\PY{l+s}{m}\PY{l+s}{o}\PY{l+s}{o}\PY{l+s}{n}\PY{l+s}{\PYZsq{}}
         \PY{l+s}{j}\PY{l+s}{o}\PY{l+s}{d}\PY{l+s}{s}\PY{l+s}{\PYZus{}}\PY{l+s}{i}\PY{l+s}{j}\PY{l+s}{o}\PY{l+s}{d}\PY{l+s}{\PYZus{}}\PY{l+s}{=}\PY{l+s}{:}\PY{l+s}{ }\PY{l+s}{\PYZsq{}}\PY{l+s}{s}\PY{l+s}{u}\PY{l+s}{n}\PY{l+s}{m}\PY{l+s}{o}\PY{l+s}{o}\PY{l+s}{n}\PY{l+s}{t}\PY{l+s}{e}\PY{l+s}{s}\PY{l+s}{t}\PY{l+s}{s}\PY{l+s}{\PYZsq{}}
         
         \PY{l+s}{N}\PY{l+s}{B}\PY{l+s}{.}\PY{l+s}{ }\PY{l+s}{p}\PY{l+s}{u}\PY{l+s}{t}\PY{l+s}{/}\PY{l+s}{x}\PY{l+s}{r}\PY{l+s}{e}\PY{l+s}{f}
         \PY{l+s}{D}\PY{l+s}{O}\PY{l+s}{C}\PY{l+s}{U}\PY{l+s}{M}\PY{l+s}{E}\PY{l+s}{N}\PY{l+s}{T}\PY{l+s}{C}\PY{l+s}{O}\PY{l+s}{M}\PY{l+s}{M}\PY{l+s}{A}\PY{l+s}{N}\PY{l+s}{D}\PY{l+s}{\PYZus{}}\PY{l+s}{z}\PY{l+s}{\PYZus{}}\PY{l+s}{ }\PY{l+s}{=}\PY{l+s}{:}\PY{l+s}{ }\PY{l+s}{\PYZsq{}}\PY{l+s}{s}\PY{l+s}{m}\PY{l+s}{o}\PY{l+s}{u}\PY{l+s}{t}\PY{l+s}{p}\PY{l+s}{u}\PY{l+s}{t}\PY{l+s}{ }\PY{l+s}{p}\PY{l+s}{r}\PY{l+s}{ }\PY{l+s}{\PYZsq{}}\PY{l+s}{\PYZsq{}}\PY{l+s}{\PYZob{}}\PY{l+s}{\PYZti{}}\PY{l+s}{N}\PY{l+s}{\PYZti{}}\PY{l+s}{\PYZcb{}}\PY{l+s}{\PYZsq{}}\PY{l+s}{\PYZsq{}}\PY{l+s}{\PYZsq{}}
         \PY{n+nl}{)}
         
         \PY{c+c1}{NB. store macro}
         \PY{l+m+mi}{4} \PY{n+nv}{put} \PY{l+s}{\PYZsq{}}\PY{l+s}{p}\PY{l+s}{r}\PY{l+s}{j}\PY{l+s}{s}\PY{l+s}{u}\PY{l+s}{n}\PY{l+s}{m}\PY{l+s}{o}\PY{l+s}{o}\PY{l+s}{n}\PY{l+s}{\PYZsq{}}\PY{o}{;}\PY{n+nv}{JSCRIPT\PYZus{}ajod\PYZus{}}\PY{o}{;}\PY{n+nv}{prjsunmoon}
\end{Verbatim}


    \begin{Verbatim}[commandchars=\\\{\}]
+-+--------------------+------+
|1|1 macro(s) put in ->|bpcopy|
+-+--------------------+------+

    \end{Verbatim}

    Once project macros are defined it's easy to configure your J session.
Open the required dictionaries and run the macor with (\texttt{rm}).

    \begin{Verbatim}[commandchars=\\\{\}]
{\color{incolor}In [{\color{incolor}22}]:} \PY{c+c1}{NB. setup project}
         \PY{n+nv}{rm} \PY{l+s}{\PYZsq{}}\PY{l+s}{p}\PY{l+s}{r}\PY{l+s}{j}\PY{l+s}{s}\PY{l+s}{u}\PY{l+s}{n}\PY{l+s}{m}\PY{l+s}{o}\PY{l+s}{o}\PY{l+s}{n}\PY{l+s}{\PYZsq{}}  \PY{o}{[} \PY{n+nv}{od} \PY{o}{;}\PY{o}{:}\PY{l+s}{\PYZsq{}}\PY{l+s}{b}\PY{l+s}{p}\PY{l+s}{c}\PY{l+s}{o}\PY{l+s}{p}\PY{l+s}{y}\PY{l+s}{ }\PY{l+s}{b}\PY{l+s}{p}\PY{l+s}{t}\PY{l+s}{e}\PY{l+s}{s}\PY{l+s}{t}\PY{l+s}{\PYZsq{}} \PY{o}{[} \PY{l+m+mi}{3} \PY{n+nv}{od} \PY{l+s}{\PYZsq{}}\PY{l+s}{\PYZsq{}}
\end{Verbatim}


    \begin{Verbatim}[commandchars=\\\{\}]
   NB. standard j scripts
   require 'debug task'
   NB. local script nouns 
   jodg\_ijod\_=: 'sunmoon'
   jods\_ijod\_=: 'sunmoontests'
   NB. put/xref
   DOCUMENTCOMMAND\_z\_ =: 'smoutput pr ''\{\textasciitilde{}N\textasciitilde{}\}'''
   

    \end{Verbatim}

    \subsubsection{\texorpdfstring{Edit your
\texttt{jodprofile.ijs}}{Edit your jodprofile.ijs}}\label{edit-your-jodprofile.ijs}

When JOD loads the script
(\texttt{\textasciitilde{}addons/general/jod/jodprofile.ijs}) is
executed. This script can be used to set up JOD the way you want. Note
how you can execute project macros when JOD loads with sentences like:

\begin{verbatim}
rm 'prjsunmoon' [ od ;:'bpcopy bptest'
\end{verbatim}

WARNING: store your profile in one of your dictionaries. This file is
reset when JAL updates the JOD addon.

    \begin{Verbatim}[commandchars=\\\{\}]
{\color{incolor}In [{\color{incolor}23}]:} \PY{c+c1}{NB. display jodprofile.ijs}
         \PY{n+nv}{smoutput} \PY{n+nv}{read\PYZus{}ajod\PYZus{}} \PY{n+nv}{jpath}\PY{l+s}{\PYZsq{}}\PY{l+s}{\PYZti{}}\PY{l+s}{a}\PY{l+s}{d}\PY{l+s}{d}\PY{l+s}{o}\PY{l+s}{n}\PY{l+s}{s}\PY{l+s}{/}\PY{l+s}{g}\PY{l+s}{e}\PY{l+s}{n}\PY{l+s}{e}\PY{l+s}{r}\PY{l+s}{a}\PY{l+s}{l}\PY{l+s}{/}\PY{l+s}{j}\PY{l+s}{o}\PY{l+s}{d}\PY{l+s}{/}\PY{l+s}{j}\PY{l+s}{o}\PY{l+s}{d}\PY{l+s}{p}\PY{l+s}{r}\PY{l+s}{o}\PY{l+s}{f}\PY{l+s}{i}\PY{l+s}{l}\PY{l+s}{e}\PY{l+s}{.}\PY{l+s}{i}\PY{l+s}{j}\PY{l+s}{s}\PY{l+s}{\PYZsq{}}
\end{Verbatim}


    \begin{Verbatim}[commandchars=\\\{\}]
NB.*jodprofile s-- JOD dictionary profile.
NB.
NB. An example JOD profile script. Save this script in
NB.
NB. \textasciitilde{}addons/general/jod/
NB.
NB. with the name jodprofile.ijs
NB.
NB. This script  is  executed  after all dictionary  objects have
NB. been created. It  can  be used  to  set up  your default  JOD
NB. working environment.
NB.
NB. WARNING: Do not dpset 'RESETME' if more than one  JOD task is
NB. active. If only one task is active RESETME's prevent annoying
NB. already open  messages that frequently result from forgetting
NB. to close dictionaries upon exiting J.
NB.
NB. Note to  J developers. A shutdown sentence (a line  of J  the
NB. interpretor  executes  before   terminating)  would  be  very
NB. useful.
NB.
NB. author: John D. Baker
NB. email:  bakerjd99@gmail.com

NB. set white space preservation on
9!:41 [ 1

NB. minimum print precision to show yyyymmdd dates (see jodage)
9!:11 [ 8

NB. set jqt windows console size - automatic for linux/mac/ios
Cwh\_j\_=: 160 24

NB. do not reset if you are running more than one JOD instance
NB. multiple JOD instances are permitted but not recommended
dpset 'RESETME'

NB. JOD interface locale - (ijod) is a good place for ad hoc JOD addons
coclass 'ijod'

NB. override EDCONSOLE - set (jconsole.exe) editor - use full path
NB. NIMP: currently EDCONSOLE is windows only - see verb (et) 
EDCONSOLE\_ajodutil\_=: '"c:\textbackslash{}Program Files\textbackslash{}Microsoft VS Code\textbackslash{}code.exe"'
NB. EDCONSOLE\_ajodutil\_=: '"c:\textbackslash{}Program Files (x86)\textbackslash{}notepad++\textbackslash{}notepad++.exe"'

NB. used by some macros: WHEREAMI=: ;0 \{ ;:'home work test'
NB. WHEREAMI=: 'home'

NB. (ijod) error/ message text
ERRIJOD00=: 'current group name (jodg\_ijod\_) not set'
ERRIJOD01=: 'current suite name (jods\_ijod\_) not set'
OKIJOD00=:  'no matches'

NB. add delete objects from current group or current suite
ag=: 3 : 0
if. wex\_ajod\_ <'jodg' do. jodg addgrp y else. jderr\_ajod\_ ERRIJOD00 end.
)
as=: 3 : 0
if. wex\_ajod\_ <'jods' do. (jods;3) addgrp y else. jderr\_ajod\_ ERRIJOD01 end.
)
dg=: 3 : 0
if. wex\_ajod\_ <'jodg' do. jodg delgrp y else. jderr\_ajod\_ ERRIJOD00 end.
)
ds=: 3 : 0
if. wex\_ajod\_ <'jods' do. (jods;3) delgrp y else. jderr\_ajod\_ ERRIJOD01 end. 
)
   
NB. referenced words not in current group
nx=: 3 : 0
if. -.wex\_ajod\_ <'jodg' do. jderr\_ajod\_ ERRIJOD00 return. end.
if. badrc\_ajod\_ gn=. grp jodg do. gn return. end.
(allrefs  \}. gn) -. gn
)
   
NB. words in current group using a word
ug=: 3 : 0
if. -.wex\_ajod\_ <'jodg' do. jderr\_ajod\_ ERRIJOD00 return. end.
if. badrc\_ajod\_ gn=. grp jodg do. gn return. end.
y usedby \}. gn
)
   
NB. generate current group and save load script
sg=: 3 : 0
if. wex\_ajod\_ <'jodg' do. mls jodg else. jderr\_ajod\_ ERRIJOD01 end. 
)

NB. top (put dictionary) words, groups, suites in revision order
tw=: revo
tg=: 2\&revo
ts=: 3\&revo

NB. run tautology as plaintest - does not stop on nonzero results
rt=: 2\&rtt

NB. run macro silently - will show explict smoutput
rs=: 1\&rm

NB. short help for groups/suites
hg=: [: hlpnl [: \}. grp
hs=: 1 hlpnl [: \}. 3 grp ]

NB. short help on put objects in revised order from code:
NB.     hr 4  NB. macro
NB.     hr 2  NB. groups
NB.  10 hr 0  NB. last ten words
hr=: 3 : 0
if. badrc\_ajod\_ w=. y revo '' do. w return. end.
y hlpnl \}. w
:
if. badrc\_ajod\_ w=. y revo '' do. w return. end.
y hlpnl x \{. \}. w
)

NB. search short help for string and list matching words
NB.     hss 'see long'   NB. search word short text 
NB.  2  hss 'see long'   NB. search group short text
NB.  4  hss 'post'       NB. search macro short text 
hss=: 3 : 0
0 hss y
:
if. badrc\_ajod\_ w=. x dnl ''   do. w return. end.
d=. x hlpnl \}. w
if. 0=\#w=. 1 >@\{ d             do. ok\_ajod\_ OKIJOD00 return. end.
if. 0=\#s=. I. (,:y) +./"1@E. w do. ok\_ajod\_ OKIJOD00 return. end.
s\&\{ \&.> d
)

NB. single line explanation 
NB.    slex 'word'         NB. word
NB.  4 slex 'jodregister'  NB. macro
NB.  1 slex 'thistest'     NB. test
slex=: 3 : 0
0 slex y
:
if. badcl\_ajod\_ sl=. x disp y do. sl return. end.
(x,8) put y;firstcomment\_\_JODtools sl
)

NB. regenerate put dictionary word cross references
reref=: 3 : 0
if. badrc\_ajod\_ n=. revo '' do. n return. end.
(n,.s) \#\textasciitilde{} -.;0\{"1 s=.0 globs\&> n=.\}.n
)

NB. handy cl doc helpers
docscr=: 3 : 'ctl\_ajod\_ (61;0;0;''NB.'') docct2\_\_UT\_\_JODobj ];.\_1 LF,y-.CR'
doctxt=: 3 : 'ctl\_ajod\_ (61;0;0;'''') docct2\_\_UT\_\_JODobj ];.\_1 LF,y-.CR'

NB. display noun on screen and return noun value
showpass=:] [ 1!:2\&2

NB. portable box drawing characters
portchars=:[: 9!:7 '+++++++++|-'"\_ [ ]

NB. windows lucida console box drawing characters
winlcchars=:[: 9!:7 (a.\{\textasciitilde{}16+i.11)"\_ [ ]

NB. edit command 
DOCUMENTCOMMAND=: 'showpass pr ''\{\textasciitilde{}N\textasciitilde{}\}'''

NB. read \& write files
read=:1!:1\&(]`<@.(32\&>@(3!:0)))
write=:1!:2 ]`<@.(32\&>@(3!:0))
readnoun=:3!:2@(1!:1\&(]`<@.(32\&>@(3!:0))))
writenoun=:([: 3!:1 [) (1!:2 ]`<@.(32\&>@(3!:0))) ]

NB. read TAB delimited table files - faster than (readtd) - see long document in (utils)
readtd2=:[: <;.\_2\&> (9\{a.) ,\&.>\textasciitilde{} [: <;.\_2 [: (] , ((10\{a.)"\_ = \{:) \}. (10\{a.)"\_) (13\{a.) -.\textasciitilde{} 1!:1\&(]`<@.(32\&>@(3!:0)))

NB. writes tables as TAB delimited LF terminated text - see long document in (utils)
writetd2=:] (1!:2 ]`<@.(32\&>@(3!:0)))\textasciitilde{} [: ([: (] , ((10\{a.)"\_ = \{:) \}. (10\{a.)"\_) [: \}.@(,@(1\&(,"1)@(-.@(*./\textbackslash{}."1@(=\&' '@])))) \# ,@((10\{a.)\&(,"1)@])) [: \}."1 [: ;"1 (9\{a.)\&,@":\&.>) [

NB. fetch edit text/macros and associate code
tt=:] ; gt
mt=:] ; 25"\_ ; gt   NB. *.txt
mj=:] ; 21"\_ ; gt   NB. *.ijs
md=:] ; 27"\_ ; gt   NB. *.markdown
mq=:] ; 30"\_ ; gt   NB. *.sql
mx=:] ; 22"\_ ; gt   NB. *.tex

NB. \textasciitilde{}user/temp object text - default j script
os=: 'ijs'\&\$: : ([: jpath '\textasciitilde{}user/temp/' , (' ' -.\textasciitilde{} ]) , '.' , ' ' -.\textasciitilde{} [)
 
NB. read text from j user temp directory
jt=:[: read os
 
NB.  load j script from j user temp
jl=: (0!:0)@jt

NB. load j script from configured j path
jp=: [: 0!:0 [: < jpath

NB. number of objects - used by various (utils) macros (sizeput, ageput, {\ldots}) if present
NOBS=: 10

NB. where am I used by some macros values 'work' or 'home'
NB. WHEREAMI=: 'work'

NB. dump drive - used by (utils) macro (dumpput) if present
NB. DUMPWINDRV=: 'h:'

NB. dump paths - used by (utils) macro (dumpput) if present
NB. DUMPPATH=: '/jod/joddumps/'
NB. DUMPPRVPATH=: '/jod/jodprvdumps/'

NB. clear dictionaries - used by (utils) macro (clearput) if present
NB. CLEARJDICS=: ;:''

NB. JOD verbs typically run from the base locale 
cocurrent 'base'

NB. examples of JOD session start ups - shows
NB. how to open dictionaries and invoke project macros

NB. set up current project (1 suppress IO, 0 or elided display)
NB. 1 rm 'prjsmughacking' [ smoutput od ;:'smugdev smug utils' [ 3 od ''
NB. 1 rm 'prjmep' [ od 'mep' [ 3 od ''


    \end{Verbatim}

    \subsubsection{Use JOD help and
documentation}\label{use-jod-help-and-documentation}

Install the (\texttt{general/joddocument}) addon to use JOD PDF help.

    \begin{Verbatim}[commandchars=\\\{\}]
{\color{incolor}In [{\color{incolor}24}]:} \PY{c+c1}{NB. location of jod.pdf \PYZhy{} install addon (general/joddocument) with JAL}
         \PY{n+nv}{smoutput} \PY{n+nv}{jpath} \PY{l+s}{\PYZsq{}}\PY{l+s}{\PYZti{}}\PY{l+s}{a}\PY{l+s}{d}\PY{l+s}{d}\PY{l+s}{o}\PY{l+s}{n}\PY{l+s}{s}\PY{l+s}{/}\PY{l+s}{g}\PY{l+s}{e}\PY{l+s}{n}\PY{l+s}{e}\PY{l+s}{r}\PY{l+s}{a}\PY{l+s}{l}\PY{l+s}{/}\PY{l+s}{j}\PY{l+s}{o}\PY{l+s}{d}\PY{l+s}{d}\PY{l+s}{o}\PY{l+s}{c}\PY{l+s}{u}\PY{l+s}{m}\PY{l+s}{e}\PY{l+s}{n}\PY{l+s}{t}\PY{l+s}{/}\PY{l+s}{p}\PY{l+s}{d}\PY{l+s}{f}\PY{l+s}{d}\PY{l+s}{o}\PY{l+s}{c}\PY{l+s}{/}\PY{l+s}{j}\PY{l+s}{o}\PY{l+s}{d}\PY{l+s}{.}\PY{l+s}{p}\PY{l+s}{d}\PY{l+s}{f}\PY{l+s}{\PYZsq{}}
         
         \PY{c+c1}{NB. opens jod.pdf if a pdf reader is available on your system}
         \PY{n+nv}{jodhelp} \PY{l+m+mi}{0}
\end{Verbatim}


    \begin{Verbatim}[commandchars=\\\{\}]
c:/j64/j64-901beta/addons/general/joddocument/pdfdoc/jod.pdf
+-+-------------------+
|1|starting PDF reader|
+-+-------------------+

    \end{Verbatim}

    \subsubsection{Summary and final
remarks}\label{summary-and-final-remarks}

\begin{enumerate}
\def\labelenumi{\arabic{enumi}.}
\item
  JOD DOES NOT BELONG IN THE J TREE
\item
  BACKUP BACKUP BACKUP
\item
  TAKE A SCRIPT DUMP
\item
  MAKE A MASTER RE-REGISTER SCRIPT
\item
  SET LIBRARY DICTIONARIES TO READONLY
\item
  KEEP REFERENCES UPDATED
\item
  DOCUMENT DICTIONARY OBJECTS
\item
  DEFINE YOUR OWN JOD SHORTCUTS
\item
  CUSTOMIZE JOD EDIT FACILITIES
\item
  DEFINE JOD PROJECT MACROS
\item
  EDIT YOUR \texttt{jodprofile.ijs}
\item
  USE JOD HELP AND DOCUMENTATION
\end{enumerate}

These are just some of the JOD practices I have found useful. As you use
JOD you will probably find your own methods.

If you find an a useful method please let me know. I can be reached at:

\begin{verbatim}
John Baker 
December 2019
bakerjd99@gmail.com
\end{verbatim}


    % Add a bibliography block to the postdoc
    
    
    
    \end{document}
